% !TEX root = ../print.tex
% Chapter 2, continued: the theory on the line, the parts of it that the plane uses. The draft's
% "The theory on the line" (its lead, lines 120-125), section 2.1 (the lattice of strata, lines
% 126-161), section 2.6 (grid pyramids, lines 260-298) and section 2.7 up to and including
% Proposition 2.23 (lines 299-316). Sections 2.2-2.5 and Proposition 2.24 are Appendix A
% (A1-line.tex). Paper II and Paper I's Remark 13.1 are unreleased and are not cited: the canonical
% discrete gauge is lem:family-dilations(3) at d = 1, and every statement the draft ported from
% Paper II is proved here or in the next chapter.
%
% The line's statements are read on the kernels. lem:line-axioms-kernels turns a family of
% operators with (A1)-(A7) and (ND) into a symmetric family in dimension 1 in the sense of
% def:family-dimension-d, and back, so that the next chapter's lemmas apply at d = 1; this is
% the draft's Remark 3.16 made into proofs.
%
% Stays prose, and why:
% - the lead "The theory on the line": positions the section in the literature, claims nothing;
% - the paragraph on knots after Lemma 2.3: reads def:knots-and-ladder at d = 1 and names the
%   phases, no claim;
% - the sentences after Definition 2.2 on self-decomposable laws: a comparison, carried by the
%   next chapter's prose and by thm:lattice-line(3);
% - the closing sentences of Example 2.5 (the preferred length, the fixed-step generator): reading;
% - the sentence after Remark 2.7 that an increment can be a lattice law: Proposition 2.21(4) is it;
% - the lead of section 2.6 and the sentence after Definition 2.20 (one increment forces the
%   ladder, one line of Lemma 2.3);
% - the readings after Proposition 2.21 (Burt and Adelson's equivalent weight run downward, the
%   rate of the geometric mask): commentary on a proved statement;
% - Remark 2.22 and the paragraph after it: a comparison with the released continuum clause
%   (V-cone v0.1, Prop. 7.5(2), the interface thm:cited-A20) and with the sampled Matern pyramid;
%   no claim beyond what Propositions 2.21 and 2.23 prove. The comparison with the released
%   clause, with the Cin ray's generator (after Example 2.5) and on locality (after
%   Proposition 2.23) is the section's one relation remark, "Relation to the continuum theory",
%   at the end of the section;
% - the lead of section 2.7, and the measurements after Proposition 2.23 (the total-variation
%   errors of the computed levels, the sampled Gaussian family's negative increments): computed
%   by scripts/grid_pyramid_exactness.py and scripts/sampled_gaussian_increments.py, reported as
%   measurements and not as proofs (WRITING.md section 4; brief, "Textbook facts");
% - the paragraph after Proposition 2.23 (the Matern family against the Gaussian on a grid): reading.

\section{The theory on the line}
\label{sec:line-theory}

An image pyramid is a sequence of images, each a smoothed and subsampled copy of the one before,
such that every level looks like every other on its own grid. The classical literature
constrains the generating kernel of one step. It requires it to be separable, normalized, symmetric
and of equal contribution [\sscite{burt1983laplacian}, p.~533], or optimal for one decimation
[\sscite{meer1987frequency}]. It does not state the property of the family that makes a pyramid a
pyramid. That property is covariance under one dilation $\lambda = q$, with $q = 2$ for the
dyadic pyramid. Under the hemigroup axioms it is the weakening of (A8) to the closed subgroup
$q^{\ZZ}$ of $(0,\infty)$. Every kernel of such a family is then a semi-selfdecomposable law in
the sense of Maejima and Naito [\sscite{maejima1998semi}]. The processes with such marginals are
the semi-selfsimilar additive processes of Maejima and Sato [\sscite{maejima1999semi}], which this
paper does not use.

This section holds the parts of the line's theory that the plane uses. The first subsection
places the stratum of one dilation between the two others, by the \Levy{} data of the kernels.
The second says which masks on a grid generate a pyramid of the axioms, and in which stratum:
the infinitely divisible masks with a logarithmic moment, in the second stratum and in no
higher one. The third computes on the grid. There the sampled Laplace kernels form an
\emph{exact cascade} on the grid: the step between any two of them is again a positive grid
kernel of unit mass, so they satisfy the cascade without approximation. The \Matern{} pyramid
computed on the grids by decimation agrees with its sampled steps at the first level only. The
word \emph{exact} has three senses on a grid, fixed in Definition~\ref{def:exact-senses}, and
each use below names its sense. The
rest of the theory on the line is Appendix~\ref{ch:line-appendix}: the gauge group, the
parametrization of one octave, the surface measure, phase entropy, the classical generating
kernels and the \Matern{} pyramid. The next chapter proves most statements of the first two
subsections again in $\RR^d$, and the proofs here read them at $d = 1$. A statement in $\RR^d$
applies at $d = 1$ through the identification $x \mapsto xe_1$ of $\RR$ with $\RR^1$, which
preserves transforms, convolution, dilations, symmetry, infinite divisibility, lattices, the
absolute value and Lebesgue measure.

Two dilations are used, as fixed in Subsection~\ref{subsec:cited-notation}. For a measure $\mu$ on
$\RR$ and $c > 0$, $T_c\mu$ is the image of $\mu$ under $x \mapsto cx$, so that
$\widehat{T_c\mu}(\omega) = \widehat\mu(c\omega)$. It is the dilation $D_c$ of measures defined
there, at $d = 1$; the next chapter, which works in $\RR^d$, writes $D_c$. For a function,
$(\mathcal{D}_\lambda f)(x) := \lambda^{-1}f(x/\lambda)$, the dilation of (A8), which preserves the
integral.

\subsection{The three strata}
\label{subsec:line-strata}

The dilations under which a family is covariant form a closed subgroup of $(0,\infty)$, and
these are of three kinds. So there are three strata: no covariance, covariance under the powers
of one ratio, and covariance under every dilation. This subsection writes each as a condition on
the \Levy{} data of the kernels.

% draft: the lead of "The theory on the line", line 124, the sentence that names a
% q^Z-covariant family; and line 132 for the canonical discrete gauge.
% CHANGED (second blind review, L4): "(A8) for the dilations in q^Z only" reads "(A8) required at the
% dilations in q^Z", so that a fully covariant family is not read as excluded. Same meaning.
\begin{definition}[families on the line covariant under one ratio]
\label{def:q-covariant-line}
\lean{Pyramids.qPowers, Pyramids.IsQCovariantLine, Pyramids.IsCanonicalDiscreteGauge,
Pyramids.IsFullyCovariantLine}\leanok
\uses{def:cascade-measurement-family}
Let $q > 1$. A family $(\Phi_{s,t})_{0 \le s \le t}$ of linear operators on $L^1(\RR)$ is
\emph{$q^{\ZZ}$-covariant} if it satisfies (A1)--(A7) and (ND) of
Definition~\ref{def:cascade-measurement-family}, and (A8) required at the dilations in $q^{\ZZ}$:
there are increasing bijections $S_\lambda$ of $[0,\infty)$, $\lambda \in q^{\ZZ}$, with
$\mathcal{D}_\lambda\Phi_{s,t} = \Phi_{S_\lambda s, S_\lambda t}\mathcal{D}_\lambda$ for all
$s \le t$. This condition is written (A8$|q^{\ZZ}$). The family is in the \emph{canonical
discrete gauge} if $S_qt = qt$ for all $t$. It is \emph{fully covariant} if it satisfies
(A1)--(A8) and (ND).
\statusT\quad\emph{(A definition; it cites nothing beyond the axioms of
Definition~\ref{def:cascade-measurement-family}. New label, from the draft's lead.)}
\end{definition}

The next lemma reads the axioms on the kernels, so that the statements of the next chapter apply
on the line.

% draft: none. Section 2.0 of the draft says that the operator and the kernel forms of the axioms
% "describe the same families" and reads the line's theory on the kernels; Remark 3.16 says that
% the next chapter's statements are the line's at d = 1. This node is the bridge both need,
% proved. New label, reported.
\begin{lemma}[the axioms on the line, read on the kernels]
\label{lem:line-axioms-kernels}
\lean{Pyramids.line_axioms_kernels_1_exists, Pyramids.line_axioms_kernels_1,
Pyramids.line_axioms_kernels_1_family, Pyramids.line_axioms_kernels_2,
Pyramids.line_axioms_kernels_2_family, Pyramids.line_axioms_kernels_3,
Pyramids.line_axioms_kernels_qcov}
\leanok
\uses{def:cascade-measurement-family, def:family-dimension-d, def:q-covariant-line,
lem:family-dilations, thm:cited-A1}
Let $(\Phi_{s,t})_{0 \le s \le t}$ be linear operators on $L^1(\RR)$.
\begin{enumerate}
\item If they satisfy (A1)--(A7), then $\Phi_{s,t}f = \mu_{s,t} * f$ for a unique family of
symmetric probability measures $\mu_{s,t}$ on $\RR$, the \emph{kernels}. These satisfy (I1) and
(I2), and (ND) holds if and only if $\mu_{s,t} \ne \delta_0$ for $s < t$. So the kernels of a
family with (A1)--(A7) and (ND) are a symmetric family in dimension $1$
(Definition~\ref{def:family-dimension-d}).
\item Conversely, if symmetric probability measures $\mu_{s,t}$ on $\RR$ satisfy (I1) and (I2),
the operators $\Phi_{s,t}f := \mu_{s,t} * f$ satisfy (A1)--(A7), and they satisfy (ND) if and
only if $\mu_{s,t} \ne \delta_0$ for $s < t$. In particular a symmetric family in dimension $1$
gives a family with (A1)--(A7) and (ND).
\item Let $\lambda > 0$ and let $S$ be an increasing bijection of $[0,\infty)$. Then
$\mathcal{D}_\lambda\Phi_{s,t} = \Phi_{Ss,St}\mathcal{D}_\lambda$ for all $s \le t$ if and only
if $T_\lambda\mu_{s,t} = \mu_{Ss,St}$ for all $s \le t$. So (A8) at the dilation $\lambda$ is
$\lambda$-covariance of the kernels, with the same bijection.
\end{enumerate}
So the $q^{\ZZ}$-covariant families (Definition~\ref{def:q-covariant-line}) are, read through their
kernels, the symmetric $q$-covariant families in dimension $1$, with the same bijections.
% CHANGED (wave 2e, 2026-11-14, a strengthening): regraded [T]. Theorem thm:cited-A1, the line
% paper's representation lemma, is itself [T] and proved from the shared core with no axiom, so
% the node spends no admitted fact; ledger A1 is off the status line.
% CHANGED (fidelity review, 2026-11-15, R26/R50/R69): the status line lists the declarations' hypotheses beside the node's (LINKAGE rule 6).
\statusT\quad\emph{(No ledger entry. The representation of each operator as the convolution with
a unique probability measure, its symmetry under (A3) and the converse are
Theorem~\ref{thm:cited-A1}, the line paper's representation lemma
\prog{V-line v0.1}{Lemma~4.1, pp.~22--23}, which the shared library proves with no axiom.
Everything else is proved here: (A6) against (I1), (A7) against (I2) through \Levy's continuity
theorem and the continuity of translation in $L^1$, (ND), and clause (3). \textbf{What it is
worth.} Every later node of this section reads the line's theory on the kernels. This node makes
the reading a statement, so that those nodes cite the next chapter at $d = 1$ rather than the
line paper's analysis. \textbf{Hypotheses} (the node's, each with what its declarations require):
(1): (A1)--(A7): same in \texttt{line\_axioms\_kernels\_1\_exists}; weaker in
\texttt{line\_axioms\_kernels\_1}, which does not require (A4) and takes the kernels $\mu_{s,t}$ as a
hypothesis, a family of probability measures representing the operators, which
\texttt{line\_axioms\_kernels\_1\_exists} supplies under (A4); (A1)--(A7) and (ND) of the last
sentence: same in \texttt{line\_axioms\_kernels\_1\_family}. (2): $\mu_{s,t}$ symmetric
probability measures: same; (I1): same; (I2): same, the continuity of
$(s,t) \mapsto \widehat\mu_{s,t}(\omega)$ on $\{0 \le s \le t\}$; a symmetric family in
dimension $1$: same in \texttt{line\_axioms\_kernels\_2\_family}. (3): $\lambda > 0$: same; $S$ an
increasing bijection of $[0,\infty)$: same; the operators of (1) or (2): weaker,
\texttt{line\_axioms\_kernels\_3} requires only that each operator be the convolution with a
probability measure $\mu_{s,t}$, not (A1)--(A7). The last sentence: (A1)--(A7) and (ND): same;
$q > 1$: same, both in \texttt{line\_axioms\_kernels\_qcov}.)}
\end{lemma}

% CHANGED (wave 2d, 2026-11-13): (A7) in (2) follows the machine-checked route, one endpoint at a
% time and the modulus of continuity of translation, in place of Skorokhod's representation; the
% last paragraph names the bijections at q^n and closes with which steps are checked.
% CHANGED (the theory paper, 2026-11-16): the closing sentence on checking and cited facts shortened; no step changed.
% CHANGED (the theory paper's blueprint round, 2026-11-16, review of Section 2, STATUS): (1) and (2)
% cite Theorem thm:cited-A1 in place of the released Lemma 4.1 it restates. No step changed.
\begin{proof}
\leanok
\uses{thm:cited-A1, lem:family-dilations}
Write $g(x) := (2\pi)^{-1/2}e^{-x^2/2}$, whose transform $\widehat g(\omega) = e^{-\omega^2/2}$
has no zero. For a finite measure $\mu$ the transform of $\mu * g$ is $\widehat\mu\,\widehat g$,
so the operator $f \mapsto \mu * f$ determines $\widehat\mu$, hence $\mu$.

(1) By Theorem~\ref{thm:cited-A1}, under (A1), (A2), (A4) and (A5) each
$\Phi_{s,t}$ is the convolution with a unique probability measure $\mu_{s,t}$, which is
symmetric under (A3). (A6) reads $\delta_0 * f = \mu_{t,t} * f$ and
$(\mu_{s,t} * \mu_{r,s}) * f = \mu_{r,t} * f$ for every $f \in L^1$, by the associativity of
convolution. By the remark above this is $\mu_{t,t} = \delta_0$ and
$\mu_{r,t} = \mu_{s,t} * \mu_{r,s}$, which is (I1). For (I2) let $(s_n,t_n) \to (s,t)$ in
$\{0 \le s \le t\}$. Since the transform of an integrable function is bounded by its $L^1$ norm,
\[
  \bigl|\widehat\mu_{s_n,t_n}(\omega) - \widehat\mu_{s,t}(\omega)\bigr|\,\widehat g(\omega)
  \le \bigl\|\mu_{s_n,t_n} * g - \mu_{s,t} * g\bigr\|_1 ,
\]
and the right side tends to $0$ by (A7). Since $\widehat g(\omega) > 0$, this is (I2). Finally
$\Phi_{s,t} = \mathrm{Id}$ if and only if $\mu_{s,t} = \delta_0$, by the remark above.

(2) The converse half of Theorem~\ref{thm:cited-A1} gives (A1)--(A5). (A6) is (I1) and the
associativity of convolution, and $\Phi_{s,t} \ne \mathrm{Id}$ if and only if
$\mu_{s,t} \ne \delta_0$, by the remark above. For (A7) write $\tau_yf := f(\cdot - y)$ and
$\Theta_f(y) := \|\tau_yf - f\|_1$, a bounded continuous function with $\Theta_f(0) = 0$, by the
continuity of translation in $L^1$. For a probability $\nu$, $\nu * f - f$ is the
$\nu$-average of $\tau_yf - f$, so $\|\nu * f - f\|_1 \le \int\Theta_f\,d\nu$. Let
$(s_n,t_n) \to (s,t)$ and move one endpoint at a time,
% CHANGED (third blind pass, Q-0767, 2026-11-15): the order of the two moves is said, so that the
% intermediate pair is always a pair s <= t.
through $(s_n, t)$ when $s_n \le t$, and through $(s, t_n)$ otherwise, where
$s \le t < s_n \le t_n$. By (I1),
$\mu_{s,t} * f - \mu_{s',t} * f = \mu_{s'',t} * (\mu_{u,v} * f - f)$ up to sign, with
$u = \min(s, s')$, $v = \max(s, s')$ and $s''$ the larger of the two, so
$\|\mu_{s,t} * f - \mu_{s',t} * f\|_1 \le \|\mu_{u,v} * f - f\|_1 \le \int\Theta_f\,d\mu_{u,v}$,
convolution with a probability being a contraction of $L^1$. In the same way
$\|\mu_{s,t} * f - \mu_{s,t'} * f\|_1 \le \int\Theta_g\,d\mu_{\min(t,t'),\max(t,t')}$ with $g$ the
convolution of $f$ with the earlier factor, and $\Theta_g \le \Theta_f$ since translation commutes
with convolution. As the two ends of an increment tend to a common limit $c$, its transform tends
to $\widehat\mu_{c,c} = 1$ by (I2), so the increment tends weakly to $\delta_0$ by \Levy's
continuity theorem, and $\int\Theta_f$ against it tends to $\Theta_f(0) = 0$. This gives (A7).

(3) For a probability $\mu$ and $f \in L^1$,
\[
  (T_\lambda\mu) * (\mathcal{D}_\lambda f)(x) = \int\lambda^{-1}f\bigl((x - \lambda y)/\lambda\bigr)\,\mu(dy)
  = \lambda^{-1}(\mu * f)(x/\lambda) = \mathcal{D}_\lambda(\mu * f)(x) .
\]
So $\mathcal{D}_\lambda\Phi_{s,t}f = (T_\lambda\mu_{s,t}) * \mathcal{D}_\lambda f$, while
$\Phi_{Ss,St}\mathcal{D}_\lambda f = \mu_{Ss,St} * \mathcal{D}_\lambda f$. Since
$\mathcal{D}_\lambda$ maps $L^1$ onto $L^1$, the two operators agree for every $f$ if and only if
the convolutions with $T_\lambda\mu_{s,t}$ and with $\mu_{Ss,St}$ agree on all of $L^1$, that is,
by the remark above, if and only if $T_\lambda\mu_{s,t} = \mu_{Ss,St}$.

For the last sentence, by (1)--(3) a family with (A1)--(A7), (ND) and (A8) at every
$\lambda \in q^{\ZZ}$ is, through its kernels, a symmetric family in dimension $1$ that is
$\lambda$-covariant at every $\lambda \in q^{\ZZ}$, with the same bijections, and conversely. In
particular it is $q$-covariant. Conversely a $q$-covariant symmetric family is
$\lambda$-covariant at every $\lambda \in q^{\ZZ}$, since the set of $\lambda$ at which it is
covariant is a group, with $S_{q^n} = S_q^n$ (Lemma~\ref{lem:family-dilations}(1)--(2)):
covariance at $\lambda$ with $S$ and at $\lambda'$ with $S'$ gives covariance at $\lambda\lambda'$
with $S \circ S'$, and at $\lambda^{-1}$ with $S^{-1}$, so $S_q$ iterated $n$ times serves $q^n$,
and its inverse iterated serves $q^{-n}$.

Every step is machine-checked, Theorem~\ref{thm:cited-A1} included.
\end{proof}

% draft: Lemma 2.1 (section 2.1), lines 128-130.
% Against the draft (moved from the status, second blind review L7): the draft proves the second
% clause from the group law of the bijections S_lambda, which it attributes to the line paper's
% Lemma 6.1. That lemma is the Fourier form of covariance; the group law is the line paper's
% Lemma 6.4(2), under the full (A8), so neither covers a set of dilations.
\begin{lemma}[closed subgroups]
\label{lem:closed-subgroups}
\lean{Pyramids.closed_subgroups_1, Pyramids.closed_subgroups_2}
\leanok
\uses{lem:line-axioms-kernels, lem:family-dilations}
The closed subgroups of $\bigl((0,\infty), \times\bigr)$ are $\{1\}$, the cyclic groups
$q^{\ZZ}$ with $q > 1$, and $(0,\infty)$. A family satisfying (A1)--(A7), (ND), and (A8) for
every dilation in a set $\Lambda \subset (0,\infty)$, satisfies (A8) for every dilation in the
closed subgroup that $\Lambda$ generates.
% CHANGED (fidelity review, 2026-11-15, R26/R50/R69): the status line lists the declarations' hypotheses beside the node's (LINKAGE rule 6).
\statusT\quad\emph{(No ledger entry. The first clause, the classification of the closed
subgroups, is proved here; the last paragraph of the proof of Lemma~\ref{lem:family-dilations}
uses the same argument. The second clause is Lemma~\ref{lem:family-dilations}(2) at $d = 1$,
read through Lemma~\ref{lem:line-axioms-kernels}; that lemma proves the group law of the
bijections and the closure from (I1), (I2) and (ND) alone. \textbf{Hypotheses} (the node's,
each with what its declarations require): first clause: a closed subgroup of
$\bigl((0,\infty),\times\bigr)$: same, \texttt{closed\_subgroups\_1} reading it as a subset of
$(0,\infty)$ containing $1$, closed under products and inverses and closed in the topology of
$(0,\infty)$. Second clause: (A1)--(A7) and (ND): same, \texttt{closed\_subgroups\_2} requires a
\texttt{CascadeCore}; $\Lambda \subset (0,\infty)$: same; (A8) at every dilation in $\Lambda$:
same.)}
\end{lemma}

% CHANGED (wave 2e, 2026-11-14): the proof adds the converse of the classification (the three
% groups are closed subgroups), which the checked proof needs and the printed proof left implicit,
% and the closing paragraph says which steps are checked; otherwise the route is the printed one.
% CHANGED (the theory paper, 2026-11-16): the closing sentence on checking and cited facts shortened; no step changed.
\begin{proof}
\leanok
\uses{lem:line-axioms-kernels, lem:family-dilations}
The logarithm is an isomorphism of $\bigl((0,\infty),\times\bigr)$ onto $(\RR,+)$ and a
homeomorphism. Let $H \ne \{0\}$ be a closed subgroup of $\RR$. Since $H = -H$, it has positive
elements; let $\ell$ be the infimum of its positive elements. Suppose $\ell > 0$. Since $H$ is
closed, $\ell \in H$, so $\ell$ is the least positive element of $H$. For $x \in H$ let $k$ be
the integer with
$k\ell \le x < (k+1)\ell$. Then $x - k\ell \in H$ lies in $[0, \ell)$, so it is $0$, and
$H = \ell\ZZ$. Suppose $\ell = 0$. Then for every $\delta > 0$ the group $H$ has an element in
$(0,\delta)$, whose multiples meet every interval of length $\delta$. So $H$ is dense and, being
closed, equal to $\RR$. With $q := e^\ell$ this gives the three groups. Conversely $\{1\}$ and
$(0,\infty)$ are closed subgroups, and so is $q^{\ZZ}$ for $q > 1$: its logarithms are the
multiples of $\log q$, a discrete and hence closed subgroup of $\RR$.

For the second clause, the kernels $\mu_{s,t}$ of the family are a symmetric family in dimension
$1$, and (A8) at a dilation $\lambda$ is $\lambda$-covariance of the kernels
(Lemma~\ref{lem:line-axioms-kernels}). By Lemma~\ref{lem:family-dilations}(2) the set $G$ of
$\lambda$ at which the kernels are $\lambda$-covariant is a closed subgroup of $(0,\infty)$. It
contains $\Lambda$, so it contains the closed subgroup that $\Lambda$ generates, and at each of
its elements Lemma~\ref{lem:line-axioms-kernels}(3) turns $\lambda$-covariance back into (A8).

Every step is machine-checked, and the cited facts enter only through
Lemma~\ref{lem:family-dilations}(2).
\end{proof}

So there are exactly three strata: no covariance, one dilation, all dilations. In the middle one,
Lemma~\ref{lem:family-dilations}(3) at $d = 1$ relabels the scales so that $S_qt = qt$, the
canonical discrete gauge. What does not survive the restriction to $q^{\ZZ}$ is the uniqueness of
that gauge (Lemma~\ref{lem:gauge-group}).

% draft: Definition 2.2 (section 2.1), line 134. The draft's sentences on what the residual is
% in a family, and on self-decomposable laws, follow as prose.
% CHANGED: the class ND_s of the draft is written out (the exponents of symmetric infinitely
% divisible laws), and the cited facts of the draft's definition are those of
% def:semi-selfdecomposable, of which this is the case d = 1 for a symmetric law.
% CHANGED (blind review, flag 18): the two facts the definition asserted (infinite divisibility,
% the folded criterion) moved to the prose after it, with their source.
% CHANGED (the theory paper's blueprint round, 2026-11-16, review of Section 2, POINTER): the
% statement points to the one definition of the notion, def:self-decomposable, in place of chapter
% 3's def:semi-selfdecomposable, and writes its identity with that definition's residual varrho_q.
% Same meaning.
% CHANGED (pass 6, the author's review of Section 2, 2026-11-19): the node no longer restates
% semi-selfdecomposability, whose one definition is def:self-decomposable (at d = 1, span q,
% residual varrho_q, as the prose before it says); it defines the folded data only, and its label
% def:semi-selfdecomposable-line becomes def:folded-data.
On the line a symmetric law $\mu$ is semi-selfdecomposable with span $q > 1$ when
$\widehat\mu(\omega) = \widehat\mu(\omega/q)\,\widehat\varrho_q(\omega)$ for some infinitely
divisible law $\varrho_q$, the residual (Definition~\ref{def:self-decomposable} at $d = 1$). The
line's statements carry such a law by its folded data.

\begin{definition}[folded data]
\label{def:folded-data}
\lean{Pyramids.IsFoldedData}\leanok
For a symmetric infinitely divisible law on $\RR$ with exponent
$\Psi(\omega) = a\omega^2 + \int_{(0,\infty)}(1 - \cos\omega v)\,\varpi(dv)$, the pair
$(a, \varpi)$ is its \emph{folded data}: $a$ is the Gaussian coefficient and $\varpi$ the image
of the \Levy{} measure under $x \mapsto |x|$.
\statusT\quad\emph{(No ledger entry is spent here. \texttt{IsFoldedData $\mu$ P} says only that
$\widehat\mu = e^{-\Psi}$ for the exponent $\Psi$ of the pair $P$; that this pair is \emph{the}
folded data of $\mu$, a function of the law, rests on the uniqueness of the \Levy{}--Khintchine
triplet, Theorem~\ref{thm:cited-A9}.)}
% CHANGED (fidelity review, 2026-11-15, R5): the status line says that IsFoldedData fixes the pair by
% the exponent identity alone, and that "the" pair rests on thm:cited-A9's uniqueness.
\end{definition}

Two facts on these laws are cited. A semi-selfdecomposable law is infinitely divisible
[\sscite{sato1999levy}, Prop.~15.5, p.~93] (Theorem~\ref{thm:cited-A43}). An infinitely divisible
law with \Levy{} measure $\nu$ is semi-selfdecomposable with span $q$ if and only if
$T_q\nu \ge \nu$, with no condition on its Gaussian part [\sscite{sato1999levy}, Thm.~15.8, p.~94]
(Theorem~\ref{thm:cited-A44}). For a symmetric law with folded data $(a, \varpi)$ the criterion
reads $T_q\varpi \ge \varpi$, since folding commutes with $T_q$ and a symmetric measure is
determined by its fold. In a
$q^{\ZZ}$-covariant family in the canonical discrete gauge the residual of a kernel
$\mu_{0,t}$ is the rescaled step $T_{q^{-1}}\mu_{t,qt}$
(Lemma~\ref{lem:discrete-covariance-levy}), infinitely divisible by
Lemma~\ref{lem:family-structure}. A self-decomposable law (Definition~\ref{def:self-decomposable})
is semi-selfdecomposable for every span. The marginals of the fully covariant families on the line
are the symmetric self-decomposable laws, as Theorem~\ref{thm:lattice-line}(3) shows in
coordinates.

% draft: Lemma 2.3 (section 2.1), lines 136-140.
% CHANGED: "In the canonical discrete gauge" is written as the hypothesis S_q t = qt, and the data
% are named as the folded data of lem:family-structure; the meaning is the draft's.
% CHANGED (blind review, flag 16): a(t) := A_t is stated.
% Against the draft (moved from the status, second blind review L7): the draft proves this from the
% line paper's Lemma 6.1, display (2.2), the uniqueness of Prop. 2.3(3) and Thm. 5.2, which carry no
% monotonicity of the data.
\begin{lemma}[discrete covariance in \Levy{} coordinates]
\label{lem:discrete-covariance-line}
\lean{Pyramids.discrete_covariance_line, Pyramids.discrete_covariance_line_semi,
Pyramids.discrete_covariance_line_annulus}
\leanok
\uses{def:q-covariant-line, lem:line-axioms-kernels, lem:family-structure, lem:discrete-covariance-levy, def:self-decomposable, def:folded-data}
Let $(\Phi_{s,t})$ satisfy (A1)--(A7) and (ND), with kernels $\mu_{s,t}$
(Lemma~\ref{lem:line-axioms-kernels}). Let $\Psi_t := -\log\widehat\mu_{0,t}$, and let
$(a(t), \varpi_t)$ be the folded data of $\mu_{0,t}$: with the \Levy{} data $(A_t, \nu_t)$ of
Lemma~\ref{lem:family-structure} at $d = 1$, $a(t) := A_t$, a number since $d = 1$, and
$\varpi_t$ is the image of $\nu_t$ under $x \mapsto |x|$, so that
$\Psi_t(\omega) = a(t)\omega^2 + \int_{(0,\infty)}(1 - \cos\omega v)\,\varpi_t(dv)$. Let $q > 1$.
Then (A8$|q^{\ZZ}$) holds with $S_qt = qt$ if and only if
\[
  \Psi_{qt}(\omega) = \Psi_t(q\omega) \ \text{ for all } t \ge 0,\ \omega \in \RR
  \qquad\Longleftrightarrow\qquad
  a(qt) = q^2a(t), \quad \varpi_{qt} = T_q\varpi_t \ \text{ for all } t > 0 .
\]
Every kernel is then semi-selfdecomposable with span $q$: $T_q\mu_{0,t} = \mu_{0,qt} =
\mu_{0,t} * \mu_{t,qt}$, and $T_q\varpi_t \ge \varpi_t$, $q^2a(t) \ge a(t)$. The family is
determined by its kernels $\mu_{0,t}$ at the scales $t$ of one \emph{fundamental annulus}
$[1, q)$.
% CHANGED (wave 2e, 2026-11-14): the assignment names the checked route, which reads the folded
% data through Lemma line-exponent-data rather than the data in R^1 of Lemma
% discrete-covariance-levy; the lemma still cites no fact of its own.
% CHANGED (fidelity review, 2026-11-15, R26/R50/R69): the status line lists the declarations' hypotheses beside the node's (LINKAGE rule 6).
\statusT\quad\emph{(No ledger entry. The proof reads Lemma~\ref{lem:line-axioms-kernels},
Lemma~\ref{lem:family-structure} at $d = 1$ for the infinite divisibility of the steps,
Lemma~\ref{lem:line-exponent-data} for the folded data, and
Lemma~\ref{lem:discrete-covariance-levy} at $d = 1$ for the fundamental annulus; it cites no
fact of its own. \textbf{Hypotheses} (the node's, each with what its declarations require):
(A1)--(A7) and (ND): same, a \texttt{CascadeCore}; the kernels $\mu_{s,t}$: same, a kernel
family of the operators; the folded data $(a(t), \varpi_t)$ of $\mu_{0,t}$: same, a pair $P(t)$
with \texttt{IsFoldedData} $(\mu_{0,t}, P(t))$ for $t \ge 0$, which exists and is unique
(Lemma~\ref{lem:line-exponent-data}(1)); \texttt{discrete\_covariance\_line\_annulus} does not require
it: weaker; $q > 1$: same. The ``then'' clauses: (A8$|q^{\ZZ}$) with $S_qt = qt$: same, required of
both families in \texttt{discrete\_covariance\_line\_annulus}, and as $S_qt = qt$ for $t \ge 0$,
the domain of $S_q$.)}
\end{lemma}

% CHANGED (wave 2e, 2026-11-14): the proof follows the machine-checked route. The two equivalences
% are proved on the line, through the transforms of the steps and the folded data of
% Lemma line-exponent-data, in place of Lemma discrete-covariance-levy in R^1 and the folding of
% its data; the monotonicity of the data comes from Lemma line-exponent-data(3); the annulus clause
% reads Lemma discrete-covariance-levy at d = 1 through the embedding of the line in R^1.
% CHANGED (the theory paper, 2026-11-16): the closing sentence on checking and cited facts shortened; no step changed.
\begin{proof}
\leanok
\uses{lem:line-axioms-kernels, lem:family-structure, lem:discrete-covariance-levy,
lem:line-exponent-data}
By Lemma~\ref{lem:line-axioms-kernels} the kernels are a symmetric family in dimension $1$, and
(A8) at $q$ with $S_qt = qt$ is $q$-covariance of the kernels with that bijection. Conversely
$q$-covariance with $S_q t = qt$ gives (A8) at every $q^n$, $n \in \ZZ$, by the last sentence of
that lemma. So (A8$|q^{\ZZ}$) with $S_qt = qt$ is $q$-covariance of the kernels with
$S_qt = qt$.

For $0 \le s \le t$, (I1) gives $\widehat\mu_{0,t} = \widehat\mu_{s,t}\,\widehat\mu_{0,s}$, and
$\widehat\mu_{0,s} = e^{-\Psi_s}$ has no zero, so
$\widehat\mu_{s,t} = e^{-(\Psi_t - \Psi_s)}$. Covariance at the pairs $(0, t)$ reads
$\widehat\mu_{0,qt}(\omega) = \widehat\mu_{0,t}(q\omega)$, that is
$\Psi_{qt}(\omega) = \Psi_t(q\omega)$. Conversely, under this condition
$\widehat\mu_{qs,qt}(\omega) = e^{-(\Psi_t(q\omega) - \Psi_s(q\omega))} = \widehat\mu_{s,t}(q\omega)$,
the transform of $T_q\mu_{s,t}$, so $T_q\mu_{s,t} = \mu_{qs,qt}$ for all $s \le t$. For the
second equivalence, by Lemma~\ref{lem:line-exponent-data}(2) the function $\Psi_t(q\,\cdot)$ has
the folded data $(q^2a(t), T_q\varpi_t)$, and by (1) of that lemma a folded pair is determined by
its exponent. So $\Psi_{qt} = \Psi_t(q\,\cdot)$ if and only if $a(qt) = q^2a(t)$ and
$\varpi_{qt} = T_q\varpi_t$. At $t = 0$ both exponents vanish, since $\mu_{0,0} = \delta_0$.

Under covariance $T_q\mu_{0,t} = \mu_{0,qt}$, and (I1) gives
$\mu_{0,qt} = \mu_{t,qt} * \mu_{0,t}$. Reading this at $q^{-1}\omega$ gives
$\widehat\mu_{0,t}(\omega) = \widehat\mu_{0,t}(q^{-1}\omega)\,\widehat\varrho_q(\omega)$ with the
residual $\varrho_q := T_{q^{-1}}\mu_{t,qt}$. The step $\mu_{t,qt}$ is infinitely divisible by
Lemma~\ref{lem:family-structure} at $d = 1$, read through the embedding $x \mapsto xe_1$ of $\RR$
in $\RR^1$, and so is its dilate; so the kernel is semi-selfdecomposable with span $q$. The data
are nondecreasing: $\Psi_{qt} - \Psi_t$ is the exponent of the symmetric infinitely divisible law
$\mu_{t,qt}$, so it lies in $\mathrm{ND}_s$ (Lemma~\ref{lem:line-exponent-data}(1)), and
Lemma~\ref{lem:line-exponent-data}(3) gives $\varpi_t \le \varpi_{qt} = T_q\varpi_t$. Also
$q^2a(t) \ge a(t)$, since $a(t) \ge 0$ and $q > 1$.

Finally, two such families whose kernels agree on $[1, q)$ have, after the embedding of the line
in $\RR^1$, symmetric families in dimension $1$ that are $q$-covariant with $S_qt = qt$ and agree
there; by Lemma~\ref{lem:discrete-covariance-levy} at $d = 1$ they agree at every pair, and so do
the families on the line, since the embedding is injective on laws. The operators are the
convolutions with their kernels, so they agree too.

Every step is machine-checked, and the cited facts enter only through
Lemmas~\ref{lem:family-structure}, \ref{lem:line-exponent-data} and
\ref{lem:discrete-covariance-levy}.
\end{proof}

\paragraph{Knots.} The knots, the knot kernels and the ladder of a $q^{\ZZ}$-covariant family in
the canonical discrete gauge are those of its kernels (Lemma~\ref{lem:line-axioms-kernels}), in
the sense of Definition~\ref{def:knots-and-ladder} at $d = 1$. The knots are the scales $q^n$,
$n \in \ZZ$, the ends of the annuli $[q^n, q^{n+1}]$. The knot kernels $\mu_{0,q^n}$ have the
exponents $\Psi_{q^n} = \Psi_1(q^n\,\cdot)$, so by Lemma~\ref{lem:discrete-covariance-line} they
are the dilates of the one knot kernel $\mu_{0,1}$. Which scale is called $1$ is a choice of gauge
(Lemma~\ref{lem:gauge-group}), and the statements below fix it. For a phase $\theta \in [1,q)$
the scales $\theta q^n$ are the \emph{knots of phase $\theta$}, which the knot-atom cascades of
Theorem~\ref{thm:knot-atom-cascades} use. The family's own knots are those of phase $1$.

% draft: Theorem 2.4 (section 2.1), lines 144-152.
% CHANGED: the last sentence of the draft's statement, that the inclusions are strict by
% Examples 2.5 and 2.6, is split off as cor:lattice-line-strict. Those examples use clause (3) of
% this theorem, so keeping the sentence here would make a cycle of \uses.
% Against the draft (moved from the status, second blind review L7): the draft cites the line
% paper's Thm. 5.2 for (1); that theorem does not state the monotonicity of the data, so it is not
% used.
\begin{theorem}[the three strata on the line]
\label{thm:lattice-line}
\lean{Pyramids.lattice_line_1, Pyramids.lattice_line_1_converse, Pyramids.lattice_line_2,
Pyramids.lattice_line_3, Pyramids.lattice_line_3_then}
\leanok
% CHANGED (wave 3d, 2026-11-14): thm:lattice-strata left the uses, since the checked proof of the
% converse in (1) builds the laws on the line; lem:line-exponent-data joined them.
\uses{lem:line-axioms-kernels, lem:family-structure, lem:line-exponent-data, lem:discrete-covariance-line, def:self-decomposable, def:folded-data, thm:cited-A4, thm:cited-A9}
Let $(\Phi_{s,t})$ satisfy (A1)--(A7) and (ND), with kernels $\mu_{s,t}$, exponents $\Psi_t$ and
folded data $(a(t), \varpi_t)$ as in Lemma~\ref{lem:discrete-covariance-line}.
\begin{enumerate}
% CHANGED: the draft's "nothing further holds in general" is stated as what it means, the
% converse: every path of folded data with these properties is the data of a family. It is
% Theorem 3.7(1)'s converse at d = 1, proved there.
% CHANGED (second blind review, M7): "Levy measures varpi_t on (0, infinity)" is written as the
% integrability it means, and "varpi_t - varpi_s a positive measure, not both zero" as
% varpi_s <= varpi_t setwise with the pairs distinct. Same meaning, stated for infinite measures.
% CHANGED (the theory paper's blueprint round, 2026-11-16): the clause is labelled "(Every
% family.)", as thm:lattice-strata(1): it holds for every family, so the strata read as nested and
% not as a partition. A label only.
\item \emph{(Every family.)} The maps $t \mapsto a(t)$ and $t \mapsto \varpi_t$ are
nondecreasing, every step $\mu_{s,t}$ is a symmetric infinitely divisible law with the folded
data $(a(t) - a(s), \varpi_t - \varpi_s)$, and $t \mapsto \Psi_t(\omega)$ is continuous for each
$\omega$. Conversely, for $t \ge 0$ let numbers $a(t) \ge 0$ and measures $\varpi_t$ on
$(0,\infty)$ with $\int\min(1, v^2)\,\varpi_t(dv) < \infty$ be given, with $a(0) = 0$ and
$\varpi_0 = 0$. Suppose that whenever $s < t$, $a(s) \le a(t)$, $\varpi_s \le \varpi_t$ setwise and
$(a(s), \varpi_s) \ne (a(t), \varpi_t)$, and that
$t \mapsto a(t)\omega^2 + \int(1 - \cos\omega v)\,\varpi_t(dv)$ is continuous for each $\omega$.
Then these are the folded data of the kernels of a family with (A1)--(A7) and (ND).
\item \emph{(Discrete covariance.)} (A8$|q^{\ZZ}$) holds with $S_qt = qt$ if and only if
$a(qt) = q^2a(t)$ and $\varpi_{qt} = T_q\varpi_t$ for all $t > 0$. Every kernel $\mu_{0,t}$ is
then semi-selfdecomposable with span $q$, and the family is a \emph{pyramid}: its kernel at the
scale $q^nt$ is the $T_{q^n}$-image of its kernel at $t$.
% CHANGED: "in the canonical gauge" is written as "after an increasing relabelling of the
% scales", and the profile form of the data is written out, as the line paper's display (2.3)
% read with v = tx; the draft's kappa = -dk is the tail measure of section 2.0 and is not needed.
% CHANGED (blind review, flag 5 and flag 2): the constant is alpha, apart from the data a'(t) of
% the relabelled family mu'; F is written out; the bijections S'_lambda t = lambda t of the
% relabelled family are part of the statement.
\item \emph{(Full covariance.)} (A8) holds if and only if there is an increasing bijection $\chi$
of $[0,\infty)$ such that the relabelled family $\mu'_{s,t} := \mu_{\chi^{-1}(s),\chi^{-1}(t)}$ has
the folded data $a'(t) = \alpha t^2$ and $\varpi'_t(dv) = k(v/t)\,v^{-1}dv$ for all $t > 0$, with
$\alpha \ge 0$, $k \ge 0$ nonincreasing on $(0,\infty)$ and
$\int_0^\infty\min(1, x^2)\,k(x)\,x^{-1}dx < \infty$. Then $\Psi'_t(\omega) = F(t\omega)$ with
$F(\omega) := \alpha\omega^2 + \int_0^\infty(1 - \cos\omega x)\,k(x)\,x^{-1}dx$, the admissible
exponent (\ref{eq:admissible-exponent}) with data $(\alpha, k)$ of Theorem~\ref{thm:cited-A4}, and the relabelled family satisfies
(A8) with the bijections $S'_\lambda t = \lambda t$, $\lambda > 0$.
\end{enumerate}
% CHANGED (wave 3d, 2026-11-14): every clause is machine-checked, so statement and proof are marked
% proved; the status line says what each citation carries on the checked route.
% CHANGED (fidelity review, 2026-11-15, R26/R50/R69): the status line lists the declarations' hypotheses beside the node's (LINKAGE rule 6).
\statusA\quad\ledger{A4}\quad\ledger{A9}\quad\emph{(\textbf{Assignment.} \ledger{A4}, the
classification \prog{V-line v0.1}{Thm.~7.3, p.~37}, carries the ``only if'' of (3): a fully
covariant family has a gauge $\chi$ and an exponent $F \not\equiv 0$ of the form of (3) with
$\widehat\mu_{s,t}(\omega) = \exp[-(F(\chi(t)\omega) - F(\chi(s)\omega))]$. \ledger{A9}, the
uniqueness of the folded \Levy{}--Khintchine pair \prog{V-line v0.1}{Prop.~2.3(3), p.~9},
identifies the data of the relabelled kernels with $(\alpha t^2, k(\cdot/t)\,v^{-1}dv)$ in (3), and
gives (ND) in the converse of (1), where equal exponents at two scales would make the two pairs
equal. Not carried by either: the direct half of (1), which is
Lemmas~\ref{lem:family-structure} and~\ref{lem:line-exponent-data}; the converse's construction,
which is Lemma~\ref{lem:line-exponent-data}(1) and (3); clause (2), which is
Lemma~\ref{lem:discrete-covariance-line}; and the ``if'' of (3), proved here. Sato's Cor.~15.11 is
not used for the density form, which is already the form (2.0) of the classification.
Theorems~\ref{thm:cited-A4} and~\ref{thm:cited-A9} are themselves theorems from \ledger{A39}, so
every clause is machine-checked. \textbf{Hypotheses} (the node's, each with what its declarations
require): (A1)--(A7) and (ND): same, a \texttt{CascadeCore}, in every declaration but
\texttt{lattice\_line\_1\_converse}; the kernels and the folded data $(a(t), \varpi_t)$: same, a
kernel family and a pair $P(t)$ with \texttt{IsFoldedData} $(\mu_{0,t}, P(t))$ for $t \ge 0$,
which exists and is unique (Lemma~\ref{lem:line-exponent-data}(1)). Converse of (1),
\texttt{lattice\_line\_1\_converse}: $a(t) \ge 0$: same; $\varpi_t$ on $(0,\infty)$ with
$\int\min(1, v^2)\,\varpi_t(dv) < \infty$: same, the fields of a \texttt{SymLevyPair};
$a(0) = 0$, $\varpi_0 = 0$: same; monotone and distinct for $s < t$: same, for $0 \le s < t$; the
continuity in $t$ for each $\omega$: same, on $t \ge 0$. (2): $q > 1$: same; $S_qt = qt$: same,
for $t \ge 0$. (3): (A8): same, in \texttt{lattice\_line\_3}; in its ``if'' and in
\texttt{lattice\_line\_3\_then}: $\chi$ an increasing bijection of $[0,\infty)$: same, its inverse
$\psi$ with \texttt{IsGauge}; $\alpha \ge 0$: same; $k \ge 0$ nonincreasing on $(0,\infty)$:
same; $\int_0^\infty\min(1, x^2)\,k(x)\,x^{-1}dx < \infty$: same, as the two halves
$\int_0^1xk(x)\,dx < \infty$ and $\int_1^\infty k(x)\,x^{-1}dx < \infty$ of an \texttt{SDProfile};
$k(0) = 0$: stronger, absent from the node, a normalisation of \texttt{SDProfile} and harmless,
since the measure $k(v/t)\,v^{-1}dv$ on $(0,\infty)$ does not read $k(0)$.)}
\end{theorem}

% CHANGED (wave 3d, 2026-11-14): the proof follows the machine-checked route. The direct half of (1)
% reads the monotonicity and the differences from Lemma line-exponent-data(3); the converse builds
% the laws on the line from Lemma line-exponent-data, instead of passing through Theorem
% lattice-strata(1) at d = 1, and (ND) is said; (3) applies Theorem cited-A4 to the operators, as
% Theorem annulus-parametrization-phase does, and its "if" is read on the relabelled operators.
% CHANGED (fourth blind pass, 2026-11-15): the folded data of mu_{0,t} in the converse of (1) by
% the uniqueness of the data; in (3), Psi_t = F(chi(t) .) from the display at s = 0; the
% integrability bound of k(./t) written out; the relabelled kernels shown to be a family with
% positive transforms before (I1) is used.
% CHANGED (the theory paper, 2026-11-16): the closing sentence on checking and cited facts shortened; no step changed.
\begin{proof}
\leanok
\uses{lem:line-axioms-kernels, lem:family-structure, lem:line-exponent-data,
lem:discrete-covariance-line, thm:cited-A4, thm:cited-A9}
(1) By Lemma~\ref{lem:line-axioms-kernels}(1) the kernels are a symmetric family in dimension
$1$. Every step $\mu_{s,t}$ is symmetric, and infinitely divisible by
Lemma~\ref{lem:family-structure} at $d = 1$, read through the embedding $x \mapsto xe_1$ of $\RR$
in $\RR^1$. By (I1) its transform is $e^{-(\Psi_t - \Psi_s)}$, so $\Psi_t - \Psi_s$ lies in
$\mathrm{ND}_s$ (Lemma~\ref{lem:line-exponent-data}(1)), and Lemma~\ref{lem:line-exponent-data}(3)
gives $a(s) \le a(t)$, $\varpi_s \le \varpi_t$, and the folded data
$(a(t) - a(s), \varpi_t - \varpi_s)$ of $\Psi_t - \Psi_s$, which are those of $\mu_{s,t}$.
Continuity of $\Psi_t(\omega) = -\log\widehat\mu_{0,t}(\omega)$ in $t$ is (I2), with
$\widehat\mu_{0,t} > 0$.

For the converse put $\Psi_t(\omega) := a(t)\omega^2 + \int(1 - \cos\omega v)\,\varpi_t(dv)$. For
$0 \le s \le t$ the hypotheses give $a(s) \le a(t)$ and $\varpi_s \le \varpi_t$, so
$\Psi_t - \Psi_s$ lies in $\mathrm{ND}_s$ (Lemma~\ref{lem:line-exponent-data}(3)), and by
Lemma~\ref{lem:line-exponent-data}(1) it is $-\log\widehat\mu_{s,t}$ for a symmetric law
$\mu_{s,t}$. These laws satisfy (I1): the transform of $\mu_{t,t}$ is $1$, that of $\delta_0$,
and $\mu_{s,t} * \mu_{r,s}$ and $\mu_{r,t}$ both have the transform $e^{-(\Psi_t - \Psi_r)}$; a law
is determined by its transform. They satisfy (I2), since
$(s,t) \mapsto e^{-(\Psi_t(\omega) - \Psi_s(\omega))}$ is continuous on $0 \le s \le t$ by the
continuity of $t \mapsto \Psi_t(\omega)$. They satisfy (ND): if $\mu_{s,t} = \delta_0$ for some
% CHANGED (2026-11-23, review round 3, supp.S13): wording.
$s < t$, then $\Psi_s = \Psi_t$, and by the uniqueness of the folded pair, Theorem~\ref{thm:cited-A9},
$(a(s), \varpi_s) = (a(t), \varpi_t)$, against the hypothesis. So they form a symmetric family in
dimension $1$, and Lemma~\ref{lem:line-axioms-kernels}(2) turns it into a family of operators with
(A1)--(A7) and (ND) whose kernels are the $\mu_{s,t}$. Since $a(0) = 0$ and $\varpi_0 = 0$,
$\Psi_0 = 0$, so $\widehat\mu_{0,t} = e^{-\Psi_t}$. Thus $(a(t), \varpi_t)$ is a folded pair of
the exponent of $\mu_{0,t}$, and by the uniqueness of the data
(Lemma~\ref{lem:line-exponent-data}(1)) it is the folded data of $\mu_{0,t}$.

(2) The equivalence and the semi-selfdecomposability are
Lemma~\ref{lem:discrete-covariance-line}, and so is $T_{q^n}\mu_{0,t} = \mu_{0,q^nt}$, by
iterating $T_q\mu_{0,t} = \mu_{0,qt}$ in both directions.

(3) Suppose (A8). Theorem~\ref{thm:cited-A4}, applied to the family of operators, whose kernels
are the $\mu_{s,t}$ since an operator determines its kernel, gives an increasing bijection $\chi$
of $[0,\infty)$ and an exponent $F \not\equiv 0$,
$F(\omega) = \alpha\omega^2 + \int_0^\infty(1 - \cos\omega x)\,k(x)\,x^{-1}dx$ with $\alpha \ge 0$,
$k \ge 0$ nonincreasing and $\int\min(1,x^2)\,k(x)\,x^{-1}dx < \infty$, such that
$\widehat\mu_{s,t}(\omega) = \exp(-[F(\chi(t)\omega) - F(\chi(s)\omega)])$. At $s = 0$, where
$\chi(0) = 0$, $\chi$ being an increasing bijection of $[0,\infty)$, and $F(0) = 0$, this is
$\Psi_t(\omega) = F(\chi(t)\omega)$ for $t \ge 0$. So the relabelled family
$\mu'_{s,t} := \mu_{\chi^{-1}(s),\chi^{-1}(t)}$ has the exponents $\Psi'_t(\omega) = F(t\omega)$.
For $t > 0$ the substitution $v = tx$, with $dx/x = dv/v$, gives
\[
  F(t\omega) = \alpha\,t^2\omega^2 + \int_0^\infty(1 - \cos\omega v)\,k(v/t)\,v^{-1}dv ,
\]
and $k(\cdot/t)$ is again nonnegative and nonincreasing. The substitution gives
\[
  \int_0^\infty\min(1, v^2)\,k(v/t)\,\frac{dv}{v} = \int_0^\infty\min(1, t^2x^2)\,k(x)\,\frac{dx}{x}
  \le \max(1, t^2)\int_0^\infty\min(1, x^2)\,k(x)\,\frac{dx}{x} < \infty ,
\]
so this is the exponent of a folded pair. By the
% CHANGED (2026-11-23, review round 3, supp.S13): wording.
uniqueness of the folded pair, Theorem~\ref{thm:cited-A9}, the data of $\mu'_{0,t}$ are $a'(t) = \alpha t^2$ and
$\varpi'_t(dv) = k(v/t)\,v^{-1}dv$.

Conversely, suppose that after a relabelling by an increasing bijection $\chi$ the data are of
this form. The relabelled kernels are again a symmetric family in dimension $1$: $\chi^{-1}$ is an
increasing bijection of $[0,\infty)$, hence an increasing homeomorphism fixing $0$, so it keeps the
cascade (I1), the continuity (I2), the inequality $s < t$ of (ND) and the laws themselves. By
Lemma~\ref{lem:family-structure} their transforms $\widehat\mu'_{0,s}$ are positive, and (I1) gives
$\widehat\mu'_{s,t} = e^{-(\Psi'_t - \Psi'_s)}$. Then $\Psi'_t(\omega) = F(t\omega)$ for $t > 0$,
by the substitution above read
backwards, and for $t = 0$, where both sides vanish since $\chi^{-1}(0) = 0$. So
$\Psi'_t(\lambda\omega) = \Psi'_{\lambda t}(\omega)$ for every $\lambda > 0$, and
\[
  \widehat{T_\lambda\mu'_{s,t}}(\omega) = \widehat\mu'_{s,t}(\lambda\omega)
  = e^{-(\Psi'_t(\lambda\omega) - \Psi'_s(\lambda\omega))}
  = e^{-(\Psi'_{\lambda t}(\omega) - \Psi'_{\lambda s}(\omega))}
  = \widehat\mu'_{\lambda s,\lambda t}(\omega) ,
\]
using $\widehat\mu'_{s,t} = e^{-(\Psi'_t - \Psi'_s)}$. So the relabelled kernels are
$\lambda$-covariant with $S'_\lambda t = \lambda t$ for every $\lambda > 0$, and by
Lemma~\ref{lem:line-axioms-kernels}(3), applied to the relabelled operators
$\Phi_{\chi^{-1}(s),\chi^{-1}(t)}$, whose kernels they are, these operators satisfy (A8) with the
bijections $S'_\lambda$. The original operators then satisfy (A8) with the increasing bijections
$S_\lambda := \chi^{-1}S'_\lambda\chi$ of $[0,\infty)$.

Every step is machine-checked, and the cited facts are \ledger{A4} and \ledger{A9} and those of
the lemmas named.
\end{proof}

Clause (1) is the first stratum, clause (2) the second and clause (3) the third. The two examples
below show them apart (Theorem~\ref{cor:lattice-line-strict}).

% draft: Example 2.5 (section 2.1), line 154. An example that asserts and argues: a proposition.
% CHANGED: "the semigroup generated by the second difference" is "by half the second difference"
% (Lindeberg 1990, eq. (28), report numbering; the librarian's reading). The kernels are written
% with the series of I_n so that the statement is self-contained, and "for no span" is stated as
% the failure of (A8 on q^Z) for every q, which is what the stratum needs.
% Moved from the status (second blind review, L7): the attribution of the kernels to Lindeberg is
% the draft's, read at the report, Thm. 4 and eq. (28); "continuous in t" is not printed there.
\begin{proposition}[Lindeberg's discrete scale space lies strictly in the first stratum]
\label{prop:lindeberg-stratum-one}
\lean{Pyramids.lindeberg_stratum_one_1, Pyramids.lindeberg_stratum_one_2,
Pyramids.lindeberg_stratum_one_3}
\leanok
% CHANGED (wave 3d, 2026-11-14): thm:cited-A39 left the uses and the status line, since the
% checked proof reads the folded data of (1) from the transform and no uniqueness of the triplet
% (the precedent of prop:stratum-one-example).
\uses{lem:line-axioms-kernels, def:family-dimension-d, lem:family-dilations,
lem:discrete-covariance-line, thm:lattice-line, thm:cited-A44}
Fix $h > 0$. For $0 \le s \le t$ let $\mu_{s,t}$ be the law on $h\ZZ$ with
$\mu_{s,t}(\{hn\}) = e^{-(t-s)}I_n(t-s)$, $n \in \ZZ$, where
$I_n(u) = \sum_{j \ge 0}(u/2)^{2j+|n|}/\bigl(j!\,(j+|n|)!\bigr)$ is the modified Bessel function of
the first kind. These are the kernels of the scale space for discrete signals of
[\sscite{lindeberg1990scalespace}], the semigroup generated on $h\ZZ$ by half the second difference.
\begin{enumerate}
\item $\widehat\mu_{s,t}(\omega) = e^{-(t-s)(1 - \cos h\omega)}$. So $\mu_{0,t}$ is compound
Poisson, with $\Psi_t(\omega) = t(1 - \cos h\omega)$, $a(t) = 0$ and $\varpi_t = t\,\delta_h$.
\item The operators $\Phi_{s,t}f := \mu_{s,t} * f$ satisfy (A1)--(A7) and (ND).
\item For no $q > 1$ and no $t > 0$ is $\mu_{0,t}$ semi-selfdecomposable with span $q$. So for no
$q > 1$ does the family satisfy (A8$|q^{\ZZ}$): it lies in the first stratum of
Theorem~\ref{thm:lattice-line} and not in the second.
\end{enumerate}
% CHANGED (wave 3d, 2026-11-14): every clause is machine-checked, so statement and proof are marked
% proved; the status line says what the citation carries on the checked route.
% CHANGED (fidelity review, 2026-11-15, R26/R50/R69): the status line lists the declarations' hypotheses beside the node's (LINKAGE rule 6).
\statusA\quad\ledger{A44}\quad\emph{(\textbf{Assignment.} \ledger{A44}, the criterion
[\sscite{sato1999levy}, Thm.~15.8, p.~94], carries the first sentence of (3): an infinitely
divisible law with a generating triplet whose \Levy{} measure is $\nu$ is semi-selfdecomposable
with span $q$ if and only if $T_q\nu \ge \nu$. The second sentence of (3) is reduced to the first,
so the criterion carries it too. Not carried by it: the Bessel series, the transform and the folded
data of (1), clause (2), and the reduction in (3), which is
Lemmas~\ref{lem:line-axioms-kernels}, \ref{lem:family-dilations} and
\ref{lem:discrete-covariance-levy}. Theorem~\ref{thm:cited-A44} is itself a theorem from
\ledger{A39} and \ledger{A42}, so every clause is machine-checked. Clause (2) uses no cited fact,
and (1) only the uniqueness of the folded data, Lemma~\ref{lem:line-exponent-data}(1), which names
the pair of the exponent the folded data.
% CHANGED (fourth blind pass, 2026-11-15): "(1) and (2) use no cited fact" corrected; the proof of
% (1) now reads the uniqueness of the folded data.
The kernels are those of Lindeberg's paper [\sscite{lindeberg1990scalespace}, Thm.~4, pp.~240--241,
and eq.~(28), p.~242]; no step
of the proof uses that attribution, and the continuity in $t$ is proved here.
\textbf{Hypotheses} (the node's, each with what its declarations require): $h > 0$: same; $0 \le s \le t$:
same, \texttt{lindeberg\_stratum\_one\_1} reading the kernel at $u = t - s \ge 0$. (3): $q > 1$:
same; $t > 0$: same; in its second sentence the family of (2): weaker,
\texttt{lindeberg\_stratum\_one\_3} requires only operators of which the $\mu_{s,t}$ are the kernels,
with any bijections, and not (A1)--(A7).)}
\end{proposition}

% CHANGED (wave 3d, 2026-11-14): the proof follows the machine-checked route. (1) reads the masses
% and the transform from the generating function of the I_n instead of the binomial masses of the
% compound Poisson law, and the folded data from the transform; (3) applies the criterion to the
% image of the kernel in R^1, and reduces the second sentence to the first through Lemmas
% family-dilations(3) and discrete-covariance-levy, where the print used Lemma
% discrete-covariance-line and the data of (1).
% CHANGED (fourth blind pass, 2026-11-15): (1) says that the kernel is symmetric and infinitely
% divisible and that its folded data are unique (lem:line-exponent-data(1), now in the uses), so
% that a(t) and varpi_t are the folded data of the law before (2) makes the kernels a family; (3)
% says that the relabelled image family is again a symmetric family.
% CHANGED (the theory paper, 2026-11-16): the closing sentence on checking and cited facts shortened; no step changed.
% CHANGED (the theory paper's blueprint round, 2026-11-16, supplement review, STATUS): the sentence
% on the checked form of (1) is dropped; the proof of (1) already closes the step with the uniqueness
% of lem:line-exponent-data(1). No step changed.
\begin{proof}
\leanok
\uses{lem:line-axioms-kernels, lem:family-structure, lem:family-dilations, lem:discrete-covariance-levy, def:self-decomposable, def:folded-data, thm:lattice-line, thm:cited-A44, lem:line-exponent-data}
(1) Let $u \ge 0$ and $\theta \in \RR$. The two exponential series give
\[
  e^{u\cos\theta} = e^{(u/2)e^{i\theta}}\,e^{(u/2)e^{-i\theta}}
  = \sum_{a, b \ge 0}\frac{(u/2)^{a+b}}{a!\,b!}\,e^{i(a-b)\theta} ,
\]
the double series converging absolutely. The pairs $(a, b)$ with $a - b = n$ are
$(j + n, j)$, $j \ge 0$, when $n \ge 0$, and $(j, j + |n|)$ when $n < 0$; in both cases
$a + b = 2j + |n|$ and $a!\,b! = j!\,(j + |n|)!$. Grouping the terms by $n$ gives
\[
  \sum_{n \in \ZZ}I_n(u)\,e^{in\theta} = e^{u\cos\theta} .
\]
At $\theta = 0$ the masses $e^{-u}I_n(u) \ge 0$ sum to $1$, so each $\mu_{s,t}$ is a probability
law. At $\theta = h\omega$, with $u = t - s$,
$\widehat\mu_{s,t}(\omega) = \sum_n e^{-u}I_n(u)\,e^{ihn\omega} = e^{-u(1 - \cos h\omega)}$. So
$\Psi_t(\omega) = t(1 - \cos h\omega) = \int_{(0,\infty)}(1 - \cos\omega v)\,t\,\delta_h(dv)$.
The transform is real and even, so $\mu_{0,t}$ is symmetric, and it is the $n$-th power of the
transform of $\mu_{0,t/n}$ for every $n$, so $\mu_{0,t}$ is infinitely divisible. Its exponent has
the folded pair $(0, t\,\delta_h)$, and the folded data of a symmetric infinitely divisible law
are unique (Lemma~\ref{lem:line-exponent-data}(1)). So $\mu_{0,t}$ has the folded data
$a(t) = 0$ and $\varpi_t = t\,\delta_h$ (Definition~\ref{def:self-decomposable}), a
finite measure, and is compound Poisson. Once (2) has made the $\mu_{s,t}$ the kernels of a family,
these are the data $(a(t), \varpi_t)$ of Theorem~\ref{thm:lattice-line}.

(2) From the transforms, $\mu_{t,t} = \delta_0$ and $\mu_{s,t} * \mu_{r,s} = \mu_{r,t}$, which is
(I1). The transform $e^{-(t-s)(1 - \cos h\omega)}$ is continuous in $(s,t)$, which is (I2). It is
even in $\omega$, so every $\mu_{s,t}$ is symmetric, and for $s < t$ it is
$e^{-2(t-s)} \ne 1$ at $\omega = \pi/h$, so $\mu_{s,t} \ne \delta_0$. So the laws form a symmetric
family in dimension $1$, and Lemma~\ref{lem:line-axioms-kernels}(2) gives (A1)--(A7) and (ND).

(3) Let $q > 1$ and $t > 0$, and read $\mu_{0,t}$ on $\RR^1$ through the embedding
$x \mapsto xe_1$. Its image $m_t$ is infinitely divisible, by (2) and
Lemma~\ref{lem:family-structure} at $d = 1$, and has the transform
$e^{-t(1 - \cos h|\omega|)} = \exp\int(e^{i\langle\omega, x\rangle} - 1)\,t\sigma(dx)$ with
$\sigma := \tfrac12(\delta_{he_1} + \delta_{-he_1})$; since $\sigma$ is symmetric,
$(0, t\sigma, 0)$ is a generating triplet of $m_t$. The measure
$T_q(t\sigma) = \tfrac{t}{2}(\delta_{qhe_1} + \delta_{-qhe_1})$ gives the sphere $|x| = h$ the mass
$0$, and $t\sigma$ gives it $t > 0$, so $T_q(t\sigma) \ge t\sigma$ fails. By the criterion
Theorem~\ref{thm:cited-A44} (where $D_q = T_q$), $m_t$ is
not semi-selfdecomposable with span $q$. Neither is $\mu_{0,t}$: the image in $\RR^1$ of an
infinitely divisible residual of $\mu_{0,t}$ at $q$ would be one of $m_t$.

Suppose now that the family satisfied (A8$|q^{\ZZ}$) for some $q > 1$, with bijections $S$. By
Lemma~\ref{lem:line-axioms-kernels}(3) the kernels are $q$-covariant with the bijection $S_q$, and
so are their images in $\RR^1$, a symmetric family in dimension $1$. By
Lemma~\ref{lem:family-dilations}(3) an increasing bijection $\chi$ of $[0,\infty)$ relabels the
image family into one that is $q$-covariant with $S_qt = qt$. The relabelled family is again a
symmetric family in dimension $1$, $\chi^{-1}$ being an increasing homeomorphism of $[0,\infty)$
that fixes $0$ and the laws being unchanged. By
Lemma~\ref{lem:discrete-covariance-levy} its kernel at the scale $1$ is semi-selfdecomposable with
span $q$. Since $\chi(0) = 0$, that kernel is $m_u$ with $u := \chi^{-1}(1) > 0$, against the first
sentence. So no $q > 1$ has (A8$|q^{\ZZ}$), and by Theorem~\ref{thm:lattice-line}(1) the family lies
in the first stratum only.

Every step is machine-checked, and the cited facts are \ledger{A44} and those of the lemmas named.
\end{proof}

The grid spacing $h$ is a preferred length, which covariance forbids. The remark at the end of the
section compares Lindeberg's generator with that of a fully covariant family.

% draft: Example 2.6 (section 2.1), lines 156-158. An example that asserts and argues: a
% proposition.
% CHANGED: the hypothesis "not all zero" is added (Example 3.9 has it); with all weights zero the
% family is trivial and fails (ND). Narrowing.
% CHANGED: the draft's last claim, that the kernels are carried by the dense group generated by
% {hq^i}, is dropped as false. Weights nonincreasing and not all zero have w_i >= w_j > 0 for all
% i <= j once some w_j > 0, so the sum over i < 0 diverges and every Levy measure nu_t is
% infinite; an infinitely divisible law with an infinite Levy measure has no atoms (Sato,
% Thm. 27.4, not verified at the page and used by no node), so it is carried by no countable set.
% The true and needed part, that no kernel is a lattice law, is kept, by
% prop:no-lattice-kernel-line. The draft's claim of a dense carrier is false whenever the weights
% are not all zero (moved from the status, second blind review L7).
% CHANGED (second blind review, M3): the weights are stated nonnegative, w_i >= 0, as
% prop:stratum-two-example states them; the proof needs it for nu_t to be a measure. Narrowing.
% CHANGED (second blind review, L2): the log-scale fraction theta is renamed eta, theta being the
% phase of the knots and vartheta the grid frequency. Rename.
% CHANGED (the theory paper's blueprint round, 2026-11-16, supplement review, NOTATION): the octave
% fraction is zeta, in the statement and the proof: eta is the Haar measure of the circle in the
% supplement, theta the phase of a knot, vartheta the grid frequency. Rename.
\begin{proposition}[the geometric-jump pyramid lies strictly in the second stratum]
\label{prop:discrete-gaussian-pyramid}
\lean{Pyramids.discrete_gaussian_pyramid_1, Pyramids.discrete_gaussian_pyramid_2,
Pyramids.discrete_gaussian_pyramid_3}
\leanok
% CHANGED (wave 3e, 2026-11-15): lem:synthesis and thm:cited-A9 left the uses. The checked proof
% builds the family by the converse of thm:lattice-line(1) and identifies no folded pair itself;
% A9 is read only inside thm:lattice-line, a proved lemma, so the node cites nothing of its own.
\uses{def:q-covariant-line, lem:line-axioms-kernels, lem:discrete-covariance-line,
thm:lattice-line, prop:no-lattice-kernel-line}
Fix $q > 1$, $h > 0$ and weights $w_i \ge 0$, $i \in \ZZ$, nonincreasing in $i$ and not all zero,
with $\sum_{i < 0}w_iq^{2i} < \infty$ and $\sum_{i \ge 0}w_i < \infty$. For $t = q^{n+\zeta}$
with $n \in \ZZ$ and $\zeta \in [0,1)$ put $w_i(\zeta) := (1-\zeta)w_i + \zeta w_{i-1}$,
\[
  \varpi_t := \sum_{i \in \ZZ}w_i(\zeta)\,\delta_{hq^{i+n}}, \qquad
  \Psi_t(\omega) := \sum_{i \in \ZZ}w_i(\zeta)\bigl(1 - \cos(hq^{i+n}\omega)\bigr),
\]
and $\varpi_0 := 0$, $\Psi_0 := 0$.
\begin{enumerate}
\item The laws $\mu_{s,t}$ with $\widehat\mu_{s,t} = e^{-(\Psi_t - \Psi_s)}$ are the kernels of a
$q^{\ZZ}$-covariant family in the canonical discrete gauge, with the folded data $a(t) = 0$ and
$\varpi_t$, and $T_q\varpi_t \ge \varpi_t$ for every $t$.
\item No increasing relabelling of the scales makes the family fully covariant: it lies in the
second stratum of Theorem~\ref{thm:lattice-line} and not in the third.
\item No kernel $\mu_{0,t}$, $t > 0$, is carried by a lattice $\xi\ZZ$, $\xi > 0$.
\end{enumerate}
% CHANGED (wave 3e, 2026-11-15): every clause is machine-checked, so statement and proof are marked
% proved. The node is regraded [T]: the uniqueness of the folded pair is read only inside
% thm:lattice-line, as for thm:surface-measure.
% CHANGED (fidelity review, 2026-11-15, R26/R50/R69): the status line lists the declarations' hypotheses beside the node's (LINKAGE rule 6).
\statusT\quad\emph{(Cites nothing of its own. The family is built by the converse of
Theorem~\ref{thm:lattice-line}(1), its covariance is Lemma~\ref{lem:discrete-covariance-line},
(2) is Theorem~\ref{thm:lattice-line}(3) against the atoms of $\varpi_t$, and (3) is
Proposition~\ref{prop:no-lattice-kernel-line}(1). The uniqueness of the folded pair and the
classification of fully covariant families are read inside Theorem~\ref{thm:lattice-line} only,
where both are theorems. Every clause is machine-checked. \textbf{Not claimed}: any statement on
the carrier of the kernels beyond the lattice clause (3). \textbf{Hypotheses} (the node's, each
with what its declarations require): $q > 1$, $h > 0$: same; $w_i \ge 0$, nonincreasing, not all
zero: same; $\sum_{i<0}w_iq^{2i} < \infty$ and $\sum_{i \ge 0}w_i < \infty$: same, as the
summability of the two families of nonnegative terms. (1): the laws $\mu_{s,t}$: same,
probability measures with these transforms for $0 \le s \le t$. (2): the operators of (1): weaker,
\texttt{discrete\_gaussian\_pyramid\_2} requires only operators of which the $\mu_{s,t}$ are the
kernels, not (A1)--(A7); the relabelling, an increasing bijection of $[0,\infty)$: same,
\texttt{IsGauge}. (3): $t > 0$, $\xi > 0$: same.)}
\end{proposition}

% CHANGED (wave 3e, 2026-11-15): the proof follows the machine-checked route. (1) builds the family
% on the line by the converse of thm:lattice-line(1), from the monotonicity of the atom weights and
% the continuity of t -> Psi_t(omega) proved here, in place of lem:synthesis at d = 1; the folded
% data are then (0, varpi_t) by construction, with no appeal to A9. (2) conjugates the relabelling
% back to the family itself before applying thm:lattice-line(3), and says why the given operators
% are those of (1).
% CHANGED (fourth blind pass, 2026-11-15): in (2), psi' is named as the inverse of the chi of
% thm:lattice-line(3).
% CHANGED (the theory paper, 2026-11-16): the closing sentence on checking and cited facts shortened; no step changed.
\begin{proof}
\leanok
\uses{lem:line-axioms-kernels, lem:discrete-covariance-line, thm:lattice-line,
prop:no-lattice-kernel-line}
(1) \emph{The pairs.} For $t > 0$ write $t = q^{n+\zeta}$ with $n \in \ZZ$ and $\zeta \in [0,1)$.
Since $w_{i-1} \ge w_i \ge 0$, the weights satisfy $w_i \le w_i(\zeta) \le w_{i-1}$. With
$1 - \cos x \le 2\min(x^2, 1)$,
\[
  w_i(\zeta)\bigl(1 - \cos(hq^{i+n}\omega)\bigr) \le 2\,w_{i-1}\min\bigl((hq^n\omega)^2q^{2i}, 1\bigr),
\]
and for any $c \in \RR$
\[
  \sum_{i \in \ZZ}w_{i-1}\min(c^2q^{2i}, 1) \le c^2q^2\sum_{j < 0}w_jq^{2j} + \sum_{j \ge 0}w_j
  < \infty ,
\]
splitting at $j = i - 1 = 0$. So the series $\Psi_t(\omega)$ converges, and, with $c = hq^n$,
$\int\min(1, v^2)\,\varpi_t(dv) < \infty$. The measure $\varpi_t$ is carried by the points
$hq^j > 0$, so $(0, \varpi_t)$ is a folded pair with the exponent $\Psi_t$; at $t = 0$ it is
$(0, 0)$.

\emph{Monotone and distinct.} The weight of $\varpi_t$ on the point $hq^j$ is
$W_t(j) := w_{j-n}(\zeta)$, the points $hq^j$ being distinct. Let $0 < s \le t$, $s = q^{m+\zeta'}$.
If $m = n$ then $\zeta' \le \zeta$ and $W_s(j) \le W_t(j)$, since $w_i(\zeta)$ is nondecreasing in
$\zeta$. If $m < n$ then $W_s(j) \le w_{j-m-1} \le w_{j-n} \le W_t(j)$. So
$\varpi_s \le \varpi_t$ setwise, and trivially so for $s = 0$. Some step of the weights is strict,
$w_{i_0} < w_{i_0-1}$: otherwise $w$ is constant, and $\sum_{i \ge 0}w_i < \infty$ makes it zero.
For $0 < s < t$ the atom $j = i_0 + m$ separates the two measures. If $m = n$ then $\zeta' < \zeta$
and $W_s(j) = w_{i_0} + \zeta'(w_{i_0-1} - w_{i_0}) < W_t(j)$. If $m < n$ then
$W_s(j) = (1-\zeta')w_{i_0} + \zeta'w_{i_0-1} < w_{i_0-1} \le w_{i_0+m-n} \le W_t(j)$, since
$\zeta' < 1$. For $s = 0 < t$, some $w_i > 0$, and $W_t(i + n) \ge w_i > 0$ while $\varpi_0 = 0$.
So $\varpi_s \ne \varpi_t$ whenever $s < t$.

\emph{Continuity.} Put $g_n(\omega) := \sum_iw_i\bigl(1 - \cos(hq^{i+n}\omega)\bigr)$. On the closed
octave $[q^n, q^{n+1}]$, with $\zeta := \log_qt - n$,
\[
  \Psi_t(\omega) = (1 - \zeta)\,g_n(\omega) + \zeta\,g_{n+1}(\omega) :
\]
for $t < q^{n+1}$ split the series along $w_i(\zeta) = (1-\zeta)w_i + \zeta w_{i-1}$ and shift the index
of the second part by one; at $t = q^{n+1}$ both sides are $g_{n+1}(\omega)$. So
$t \mapsto \Psi_t(\omega)$ is continuous on each closed octave, hence on $(0,\infty)$. Further
$g_n(\omega) = G(q^n)$ with $G(c) := \sum_iw_i\bigl(1 - \cos(hq^ic\omega)\bigr)$, a series of
continuous functions dominated on $[-1,1]$ by the convergent series
$2\sum_iw_{i-1}\min((h\omega)^2q^{2i}, 1)$, so $G$ is continuous there, with $G(0) = 0$. Hence
$g_n(\omega) \to 0$ as $n \to -\infty$, and since $0 \le \Psi_t(\omega) \le \max(g_n(\omega),
g_{n+1}(\omega))$ on the octave of $t$, $\Psi_t(\omega) \to 0 = \Psi_0(\omega)$ as $t \downarrow 0$.

\emph{The family.} The pairs $(0, \varpi_t)$ thus satisfy the hypotheses of the converse in
Theorem~\ref{thm:lattice-line}(1), which gives a family with (A1)--(A7) and (ND) whose kernels
$\mu^0_{s,t}$ have the folded data $(0, \varpi_t)$. By (I1) their transforms are
$\widehat\mu^0_{s,t} = e^{-(\Psi_t - \Psi_s)}$, so laws with these transforms exist, and any such laws
are the $\mu^0_{s,t}$, a law being determined by its transform. They are therefore the kernels of
that family, with the folded data $(0, \varpi_t)$. For $t > 0$, $qt = q^{(n+1)+\zeta}$, so
$\Psi_{qt}(\omega) = \Psi_t(q\omega)$, and $\Psi_0 = 0$. By Lemma~\ref{lem:discrete-covariance-line}
the family satisfies (A8$|q^{\ZZ}$) with $S_qt = qt$, so it is $q^{\ZZ}$-covariant in the canonical
discrete gauge (Lemma~\ref{lem:line-axioms-kernels}), and $T_q\varpi_t \ge \varpi_t$. Directly:
$T_q\varpi_t$ gives the point $hq^{i+n}$ the weight $w_{i-1}(\zeta) \ge w_i(\zeta)$.

(2) Let operators $\Phi_{s,t}$ have the kernels $\mu_{s,t}$. An operator with a kernel is
convolution by it, so they are the operators of (1). Suppose an increasing bijection $\psi$ of
$[0,\infty)$ made the relabelled operators $\Phi_{\psi(s),\psi(t)}$ satisfy (A8), with bijections
$S_\lambda$, and let $\chi := \psi^{-1}$. Then the family itself satisfies (A8), with the increasing
bijections $\psi \circ S_\lambda \circ \chi$. By Theorem~\ref{thm:lattice-line}(3) there are an
increasing bijection $\chi'$ of $[0,\infty)$ and data $(\alpha, k)$ such that the relabelled family
$\mu_{\psi'(s),\psi'(t)}$, $\psi' := \chi'^{-1}$, has the folded data $\varpi_{\psi'(t)}(dv) =
k(v/t)\,v^{-1}dv$ for $t > 0$, a measure with a density, which gives every point the mass $0$. But
$u := \psi'(1) > 0$, and $\varpi_u$ gives some point $hq^j$ a positive mass, as shown in (1). So no
such $\psi$ exists.

(3) The family is $q^{\ZZ}$-covariant, and this is
Proposition~\ref{prop:no-lattice-kernel-line}(1).

Every step is machine-checked, and the cited facts enter only through
Theorem~\ref{thm:lattice-line}.
\end{proof}

The family repairs Lindeberg's discrete scale space for covariance. Lindeberg's kernels, his
discrete analogue of the Gaussian, have jumps of the one size $h$, and here each octave adds jumps
of the next size $hq^i$. That is the reason for the name; the family has no Gaussian part. Its
\Levy{} measures are atomic, and in the plane no rotation-invariant law has a
marginal with an atomic \Levy{} measure (Proposition~\ref{prop:what-does-not-lift}), which is why
the next chapter's examples put their jumps on spheres.

% draft: Theorem 2.4 (section 2.1), its last sentence, line 150, with its proof's "Nesting is
% restriction" (line 152).
% CHANGED: split from Theorem 2.4 (thm:lattice-line), so that the examples it rests on can use
% that theorem without a cycle of \uses.
\begin{theorem}[the three strata are distinct]
\label{cor:lattice-line-strict}
\lean{Pyramids.lattice_line_strict_imp, Pyramids.lattice_line_strict_lower,
Pyramids.lattice_line_strict_upper, Pyramids.lattice_line_strict}
\leanok
\uses{thm:lattice-line, prop:lindeberg-stratum-one, prop:discrete-gaussian-pyramid}
Fix $q > 1$. As conditions on families with (A1)--(A7) and (ND), full covariance implies
(A8$|q^{\ZZ}$), and both implications of the chain
\emph{fully covariant} $\Rightarrow$ \emph{$q^{\ZZ}$-covariant} $\Rightarrow$ \emph{(A1)--(A7) and
(ND)} are strict: Lindeberg's discrete scale space
(Proposition~\ref{prop:lindeberg-stratum-one}) is in the first stratum of
Theorem~\ref{thm:lattice-line} and not in the second, and the geometric-jump pyramid
(Proposition~\ref{prop:discrete-gaussian-pyramid}) is in the second and not in the third.
% CHANGED (wave 3e, 2026-11-15): every clause is machine-checked, so statement and proof are marked
% proved.
% CHANGED (fidelity review, 2026-11-15, R26/R50/R69): the status line lists the declarations' hypotheses beside the node's (LINKAGE rule 6).
\statusT\quad\emph{(Cites nothing beyond the two examples. Every clause is machine-checked.
\textbf{Hypotheses} (the node's, each with what its declarations require): $q > 1$: same, in all four;
families with (A1)--(A7) and (ND): same, a \texttt{CascadeCore}; full covariance in the
implication: same, \texttt{IsFullyCovariantLine} in \texttt{lattice\_line\_strict\_imp}. The
instances: $h > 0$ for Lindeberg's scale space: same, in \texttt{lattice\_line\_strict\_lower};
the weights of Proposition~\ref{prop:discrete-gaussian-pyramid}: same, in
\texttt{lattice\_line\_strict\_upper}; \texttt{lattice\_line\_strict} requires only $q > 1$, and fixes
the instances.)}
\end{theorem}

% CHANGED (wave 3e, 2026-11-15): the proof follows the machine-checked route; it names the instances
% of the two examples that witness the strictness for a given q, as thm:lattice-strata-strict does.
\begin{proof}
\leanok
\uses{prop:lindeberg-stratum-one, prop:discrete-gaussian-pyramid}
(A8) requires the covariance at every $\lambda > 0$, so in particular at every $\lambda \in q^{\ZZ}$,
which is (A8$|q^{\ZZ}$). Lindeberg's discrete scale space, at any $h > 0$, satisfies (A1)--(A7)
and (ND) and not (A8$|q^{\ZZ}$) (Proposition~\ref{prop:lindeberg-stratum-one}(2)--(3)). The
geometric-jump pyramid satisfies (A8$|q^{\ZZ}$), and no relabelling of its scales makes it fully
covariant (Proposition~\ref{prop:discrete-gaussian-pyramid}(1)--(2)). It is not fully covariant in
its own labelling either, the identity being a relabelling. For the given $q$ one instance is
$h = 1$ with the weights $w_i = 1$ for $i < 0$ and $w_i = 0$ for $i \ge 0$: they are nonnegative,
nonincreasing and not all zero, $\sum_{i<0}q^{2i} < \infty$, and $\sum_{i \ge 0}w_i = 0$.
\end{proof}

% draft: Remark 2.7 (section 2.1), line 160. A remark that asserts and argues: a proposition.
% CHANGED: narrowed as the draft's correction (4) of its checks narrowed it, to the kernels from
% scale zero; the argument is Proposition 3.12's at d = 1. The draft's reading of the line paper's
% Lemma 6.2 (the check R5) is not used: that lemma assumes the full (A8) and its proof needs a
% ratio with neither it nor its inverse an integer, so it does not cover q^Z with q = 2 (the
% librarian, records/pyramids/VERBATIM-V-line.md, A10). The sampling sentence is clause (2), with
% "a continuous density" so that sampling is defined.
\begin{proposition}[no kernel from scale zero is a lattice law, and the pyramid principle]
\label{prop:no-lattice-kernel-line}
\lean{Pyramids.no_lattice_kernel_line_1, Pyramids.no_lattice_kernel_line_2}
\leanok
\uses{def:q-covariant-line, lem:line-axioms-kernels, prop:no-lattice-kernel,
lem:discrete-covariance-line}
Let $q > 1$ and let $(\Phi_{s,t})$ be $q^{\ZZ}$-covariant, with kernels $\mu_{s,t}$.
\begin{enumerate}
\item No kernel $\mu_{0,t}$ with $t > 0$ is carried by a lattice $\xi\ZZ$, $\xi > 0$.
\item Suppose $S_qt = qt$, and that $\mu_{0,t}$ has a continuous density $\phi_t$ for some
$t > 0$. Then $\mu_{0,q^nt}$ has the continuous density $\phi_{q^nt}(x) = q^{-n}\phi_t(q^{-n}x)$
for every $n \in \ZZ$, so $q^n\phi_{q^nt}(q^nhm) = \phi_t(hm)$ for every $h > 0$ and $m \in \ZZ$:
sampled at the scale $q^nt$ on the grid $q^nh\ZZ$, the kernel is the same array at every level,
up to the factor $q^n$.
\end{enumerate}
% CHANGED (fidelity review, 2026-11-15, R26/R50/R69): the status line lists the declarations' hypotheses beside the node's (LINKAGE rule 6).
\statusT\quad\emph{(No ledger entry. With covariance under every dilation, clause (1) is the line
paper's Lemma~6.2 (p.~29), which is not used: its hypothesis is the full (A8).
\textbf{Hypotheses} (the node's, each with what its declarations require): $q > 1$ and
$q^{\ZZ}$-covariance: same, \texttt{IsQCovariantLine}, which holds $q > 1$, (A1)--(A7), (ND) and
(A8$|q^{\ZZ}$); the kernels $\mu_{s,t}$: same, a kernel family. (1): $t > 0$, $\xi > 0$: same.
(2): $S_qt = qt$: same, for $t \ge 0$; $t > 0$: same; a continuous density $\phi_t$: same, a
continuous $\phi \ge 0$ with $\mu_{0,t} = \phi\,dx$; $h > 0$, $m \in \ZZ$: same.)}
\end{proposition}

% CHANGED (wave 2f, 2026-11-14): the proof follows the machine-checked route. It says how the
% family on the line is read in R^1, and why the continuous density is unique.
\begin{proof}
\leanok
\uses{lem:line-axioms-kernels, prop:no-lattice-kernel, lem:discrete-covariance-line,
lem:family-dilations, def:family-dimension-d}
(1) By Lemma~\ref{lem:line-axioms-kernels} the kernels form a symmetric family in dimension $1$
that is $q$-covariant. Its image in $\RR^1$ under $x \mapsto xe_1$ is a symmetric $q$-covariant
family in the sense of Definition~\ref{def:family-dimension-d}, and the lattice $\xi\ZZ$ goes to
the subgroup generated by $\xi e_1$, which is discrete: its nonzero points have norm at least
$\xi$. Proposition~\ref{prop:no-lattice-kernel} at $d = 1$ says that no kernel of the image, at
$t > 0$, is carried by it, and the image of $\mu_{0,t}$ is carried by it exactly when $\mu_{0,t}$ is
carried by $\xi\ZZ$.

(2) In the canonical gauge, $T_q\mu_{0,t} = \mu_{0,qt}$ for every $t \ge 0$
(Lemma~\ref{lem:discrete-covariance-line}), and iterating this and its inverse,
$\mu_{0,q^nt} = T_{q^n}\mu_{0,t}$ for every $n \in \ZZ$. If $\mu$ has the density $\phi$, then
$T_c\mu$ has the density $c^{-1}\phi(x/c)$, by the substitution $x = cy$. This density is
continuous. Two continuous densities of one law agree almost everywhere, hence everywhere, since a
nonempty open set has positive Lebesgue measure. So it is the continuous density of
$\mu_{0,q^nt}$, and at $x = q^nhm$ it is $q^{-n}\phi_t(hm)$.
\end{proof}

A step $\mu_{s,t}$ with $s > 0$ can be a lattice law, and that is what a grid pyramid is
(Definition~\ref{def:grid-pyramid-line}, Proposition~\ref{prop:grid-pyramids-infinitely-divisible}(4)).
The pyramid of the axioms is a family of continuum kernels. The grids of an implemented pyramid
are the sampling of Subsection~\ref{subsec:sampled-laplace}, not part of the family. What
covariance says about the grids is clause (2), the pyramid principle.

\subsection{Grid pyramids: which masks have a pyramid, and in which stratum}
\label{subsec:grid-pyramids-line}

Proposition~\ref{prop:classical-kernels-no-exact-cascade} excludes the masks of bounded support.
A mask need not be finite to be implementable, so bounded support does not settle the question the
pyramid literature poses: given a positive mask on a grid, is there a family of the axioms whose
ladder it generates? This subsection answers it on a grid. The continuum form of the question is
answered in \prog{V-cone v0.1}{Prop.~7.5(2), p.~65} (Theorem~\ref{thm:cited-A20}), and the lattice
form is a different statement, not a corollary of it. Throughout, $h > 0$ is the grid spacing and $\vartheta = h\omega$ the grid
frequency.

% draft: Definition 2.20 (section 2.6), line 264. Its last sentence (one increment forces the
% ladder) is a claim, and follows as prose with its one-line reason.
% CHANGED: "admissible" and "fully admissible" are stated in the canonical discrete gauge, since
% the increments named sit at the scales q^n; the draft leaves the gauge implicit.
% CHANGED (blind review, flag 5; the session's convention): the draft's "admissible" and "fully
% admissible" grid pyramid are renamed "exact" (Definition 5.2's word) and "exact in the third
% stratum", so that "admissible" means only the admissible exponents of section 2.0.
\begin{definition}[the grid pyramid of a mask]
\label{def:grid-pyramid-line}
\lean{Pyramids.IsGridMask, Pyramids.gridLadderStep, Pyramids.CarriesGridLadder,
Pyramids.IsExactGridPyramid, Pyramids.IsExactGridPyramidThird}\leanok
\uses{def:q-covariant-line, lem:line-axioms-kernels, def:knots-and-ladder}
Let $q > 1$ and $h > 0$, and let $w \ne \delta_0$ be a symmetric probability law on $h\ZZ$, the
\emph{mask}. The \emph{grid pyramid} of $(w, q)$ is the ladder $\mu_{q^n,q^{n+1}} := T_{q^n}w$,
$n \in \ZZ$, read in the canonical discrete gauge $S_qt = qt$, so that its knots are the scales
$q^n$: the same array, applied on the grid $q^nh\ZZ$ at every level. It is
\emph{exact} if some $q^{\ZZ}$-covariant family in the canonical discrete gauge has these
steps, through its kernels (Lemma~\ref{lem:line-axioms-kernels}), and \emph{exact in the third stratum} if
some such family is fully covariant.
\statusT\quad\emph{(A definition; it cites nothing. It is the operator form, at $d = 1$, of
Definition~\ref{def:exact-grid-pyramid}, which is stated on laws: a hemigroup of laws with weak
continuity and covariance at $q$, without (ND). The two are not identified by definition. Their
equivalence is a statement, Proposition~\ref{prop:line-and-arrays}(1). The word ``admissible'' is
kept for exponents, (\ref{eq:admissible-exponent}).)}
\end{definition}

This is the pyramid principle read on the steps rather than on the kernels. Under
(A8$|q^{\ZZ}$) with $S_qt = qt$ the one requirement $\mu_{1,q} = w$ already forces the ladder,
since $T_{q^n}\mu_{1,q} = \mu_{q^n,q^{n+1}}$ (Lemma~\ref{lem:discrete-covariance-line}). So either
notion is a condition on one step. Definition~\ref{def:exact-grid-pyramid} is the notion in
$\RR^d$, stated on laws; this definition is its operator form on the line. That the two agree at
$d = 1$ is Proposition~\ref{prop:line-and-arrays}(1).

% draft: Proposition 2.21 (section 2.6), lines 266-291.
% CHANGED: the transform of a mask is the transform of section 2.0, w-hat(omega) for w as a law on
% R, in place of the draft's grid transfer at h omega; (2.10) reads accordingly. Notation only.
% CHANGED: renamed lambda_a to lambda^L_a, the line mask of ratio a (BLUEPRINT-INPUTS), and
% c = 2a/(1-a)^2 to 2 varsigma, varsigma := a/(1-a)^2 (BLUEPRINT-INPUTS: varsigma for a/(1-a)^2;
% c is the Matern range factor of section 2.7).
% CHANGED: in (2), the sentence that the families carrying the ladder "form its modulation orbit
% (section 2.4)" and the reading as Burt and Adelson's equivalent weighting function run downward
% go to the prose after the proposition; they name and read, and claim nothing the statement
% does not.
% CHANGED: in (2), "carried by hZ[1/q]_{>0}" is written as the union of the q^{-m}hZ_{>0}, which
% is what the proof gives for every real q > 1.
% CHANGED (blind review, flags 5 and 6): "admissible" / "fully admissible" grid pyramid are
% "exact" / "exact in the third stratum" (def:grid-pyramid-line); the lattice is xi Z, the letter
% c staying the Matern range factor. The displays the draft numbers (2.9), (2.10), (2.11) carry
% the labels eq:grid-pyramid-criterion, eq:grid-pyramid-knot-kernel,
% eq:geometric-pyramid-knot-kernel.
% CHANGED (second blind review, M1): the clause on an infinite logarithmic moment names its
% hypothesis, "w infinitely divisible and int log_+ v varpi_w(dv) = infinity", which is the
% hypothesis of prop:ladder-knot-kernel(3) that its proof uses. Narrowing.
% CHANGED (second blind review, M2): "the B-spline masks" is narrowed to those of odd length. An
% even-length B-spline mask, such as [1,3,3,1]/8, is symmetric about a half-integer and is not a
% mask in the sense of def:grid-pyramid-line. Narrowing.
% Moved from the status (second blind review, L7): the draft's own proof cites Sato's Cor. 24.6,
% 24.5, 24.4, Thm. 25.3 with Prop. 25.4(iii), the line paper's Thm. 7.3 with (2.3), and Schilling,
% Song and Vondracek's Example 13.13, and proves (4) by a Bernstein-function argument and Jensen's
% formula.
\begin{proposition}[grid pyramids and infinitely divisible masks]
\label{prop:grid-pyramids-infinitely-divisible}
\lean{Pyramids.grid_pyramids_id_1, Pyramids.grid_pyramids_id_2, Pyramids.grid_pyramids_id_2_knot,
Pyramids.grid_pyramids_id_2_carrier, Pyramids.grid_pyramids_id_2_no_kernel,
Pyramids.grid_pyramids_id_2_not_every, Pyramids.grid_pyramids_id_2_coupling,
Pyramids.grid_pyramids_id_3, Pyramids.grid_pyramids_id_3_general, Pyramids.grid_pyramids_id_4,
Pyramids.grid_pyramids_id_4_binomial, Pyramids.grid_pyramids_id_4_bounded}
\leanok
% CHANGED (wave 3e, 2026-11-15): lem:line-exponent-data joins the uses; the checked proof reads the
% folded data of an infinitely divisible mask, and their uniqueness, from its clause (1).
\uses{def:grid-pyramid-line, def:knots-and-ladder, lem:line-axioms-kernels, lem:lattice-infinitely-divisible,
prop:no-bounded-step, prop:ladder-knot-kernel, lem:discrete-covariance-line, lem:family-dilations,
thm:lattice-line, prop:no-lattice-kernel-line, prop:ladder-not-every-frequency,
lem:line-exponent-data}
Let $q > 1$, $h > 0$, and let $w \ne \delta_0$ be a symmetric probability law on $h\ZZ$.
\begin{enumerate}
\item \emph{(Lattice carrier.)} The mask $w$ is infinitely divisible if and only if it is compound
Poisson,
\[
  \widehat w(\omega) = \exp\Bigl(-\int_{(0,\infty)}(1 - \cos\omega v)\,\varpi_w(dv)\Bigr),
\]
with a finite measure $\varpi_w \ne 0$ carried by $h\ZZ_{>0}$, its folded \Levy{} measure. Then
$\widehat w > 0$ everywhere and the support of $w$ is unbounded. So a mask of bounded support is
not infinitely divisible.
\item \emph{(Criterion, and the covariant kernel.)} The grid pyramid of $(w, q)$ is exact if
and only if
\begin{equation}
\label{eq:grid-pyramid-criterion}
  w \text{ is infinitely divisible and } \int_{(0,\infty)}\log_+v\;\varpi_w(dv) < \infty ,
\end{equation}
equivalently if and only if $w$ is infinitely divisible and $\int\log_+|x|\,w(dx) < \infty$. Its
knot kernel is then unique: every family carrying the ladder has $a(1) = 0$ and
\begin{equation}
\label{eq:grid-pyramid-knot-kernel}
  \varpi_1 = \sum_{m \ge 1}T_{q^{-m}}\varpi_w, \qquad
  \widehat\mu_{0,1}(\omega) = \prod_{m \ge 1}\widehat w(q^{-m}\omega),
\end{equation}
the product being the limit of the transforms of $T_{q^{-1}}w * \dots * T_{q^{-M}}w$ as
$M \to \infty$. The measure $\varpi_1$ is carried by the union of the sets $q^{-m}h\ZZ_{>0}$,
$m \ge 1$, and by no lattice $\xi\ZZ$. When $w$ is infinitely divisible and
$\int\log_+v\;\varpi_w(dv) = \infty$, the partial products
have no limit that is a characteristic function, and they tend to $0$ at almost every $\omega$,
% CHANGED (skeleton review, m8): the parenthesis names the ratio at which "not in general" is
% shown, the typed form; under the fixed q of the clause the sentence read as a claim at that q.
though not in general at every $\omega \ne 0$ (at $q = 2$, by
Proposition~\ref{prop:ladder-not-every-frequency}): there is no kernel. The families carrying the ladder
are not unique. One of them is the \emph{uniform coupling}, with $a(t) = 0$ and
$\varpi_t = \varpi_1 + (\log_qt)\,\varpi_w$ for $t \in [1,q]$, each atom of the mask switched on
linearly in $\log t$ across the octave.
% CHANGED (second blind review, L6): "Equivalently" reads "More generally". The proof shows that
% the second sentence implies the first, and the converse argument is dropped as adding nothing
% the statement needs. Both sentences are still asserted; the claimed link is the weaker one.
\item \emph{(No grid pyramid is exact in the third stratum.)} For no $q > 1$, $h > 0$ and mask $w$ is the
grid pyramid of $(w, q)$ exact in the third stratum. More generally, a symmetric law $\varrho \ne \delta_0$
carried by a lattice $\xi\ZZ$, $\xi > 0$, is a step $\mu_{s,t}$ of no fully covariant family. So
every family carrying an exact grid pyramid lies in the second stratum of
Theorem~\ref{thm:lattice-line} and not in the third: the pyramid has a covariant kernel, unique at
the knots, and no fully covariant family interpolates its levels.
\item \emph{(Instances.)} For $0 < a < 1$ the \emph{sampled Laplace mask}
$\lambda^{\mathrm{L}}_a$, the law on $h\ZZ$ with $\lambda^{\mathrm{L}}_a(\{hn\}) =
\tfrac{1-a}{1+a}a^{|n|}$, is infinitely divisible, with all moments finite. So the
\emph{geometric pyramid}, the grid pyramid of $(\lambda^{\mathrm{L}}_a, q)$, is exact for
every $q > 1$, with the unique knot kernel
\begin{equation}
\label{eq:geometric-pyramid-knot-kernel}
  \widehat\mu_{0,1}(\omega) = \prod_{m \ge 1}\Bigl[1 + 2\varsigma\bigl(1 - \cos(q^{-m}h\omega)\bigr)\Bigr]^{-1},
  \qquad \varsigma := \frac{a}{(1-a)^2},
\end{equation}
and by (3) it lies in no stratum above the second.
% CHANGED (fourth blind pass, 2026-11-15, a strengthening to the checked declaration
% grid_pyramids_id_4_bounded, which reads every mask of bounded support): the sentence was "The
% binomial mask \tfrac1{16}[1,4,6,4,1] and the B-spline masks of odd length have bounded support,
% so they have no exact grid pyramid."; "B-spline mask" was nowhere defined.
No mask of bounded support has an exact grid pyramid. Among these masks are two families of odd
length: the binomial masks with the weights $4^{-k}\binom{2k}{k+n}$ at $hn$, $|n| \le k$,
$k \ge 1$, which are the refinement masks of the B-splines of odd degree $2k - 1$ and include
$\tfrac1{16}[1,4,6,4,1]$ at $k = 2$, and the sampled B-spline masks, the values $\beta(n)$, $n \in \ZZ$, of a B-spline $\beta$
of degree at least $2$ with unit knot spacing, centred at $0$, placed at $hn$ and renormalized to unit
mass.
\end{enumerate}
% CHANGED (wave 3e, 2026-11-15): every clause is machine-checked, so statement and proof are marked
% proved; the status line says how the line is carried to R^1 and which nodes are read there.
% CHANGED (2026-11-15, a strengthening): unused hypotheses dropped from the Lean of four
% declarations, listed at the end of the status line; the print keeps the node's setting.
% CHANGED (fidelity review, 2026-11-15, R26/R50/R69): the status line lists the declarations' hypotheses beside the node's (LINKAGE rule 6).
% CHANGED (fidelity review, 2026-11-15, R64): the status line says which named masks of (4) have a declaration; "Every clause is machine-checked" is qualified accordingly.
\statusT\quad\emph{(No ledger entry: every cited fact is carried by the nodes used, at $d = 1$:
Lemma~\ref{lem:lattice-infinitely-divisible} for the lattice carrier,
Proposition~\ref{prop:no-bounded-step} for the unbounded support, and
Proposition~\ref{prop:ladder-knot-kernel}(1)--(4)(a) for (2)--(3), which the next chapter proves
in $\RR^d$, and Proposition~\ref{prop:ladder-not-every-frequency} for the frequency at which the
partial products of (2) do not tend to $0$; Lemma~\ref{lem:line-exponent-data}(1) gives the
folded data of an infinitely divisible mask and their uniqueness. Each of these nodes is read on
the image of the line in $\RR^1$ under $x \mapsto xe_1$. Sato's Cor.~24.4--24.6 and Thm.~25.3
are spent inside those nodes, or not at all. Clause (4) is proved from the logarithm series of the
two one-sided geometric factors of the transform, which gives the mask's \Levy{} measure in closed
form. Every clause is machine-checked, with the one exception in (4) named next. Of the masks
named in (4), only $\tfrac1{16}[1,4,6,4,1]$ has a declaration of its own,
\texttt{grid\_pyramids\_id\_4\_binomial} (the mask \texttt{binomialMask5}); the binomial masks at
$k \ne 2$ and the sampled B-splines $\beta(n)$ are not defined in the Lean, and reach
\texttt{grid\_pyramids\_id\_4\_bounded}, which holds for every mask of bounded support, only
through the sentence of the proof that they are masks of bounded support, which is not
machine-checked. \textbf{Not claimed}: anything about the masks' equivalent weights beyond the
product formula. \textbf{Hypotheses} (the node's, each with what its declarations require; four
declarations had unused hypotheses dropped on 2026-11-15): $q > 1$: same where the clause reads
$q$; \texttt{grid\_pyramids\_id\_1} and \texttt{grid\_pyramids\_id\_3\_general} do not require it, their
sentences not reading $q$, and \texttt{grid\_pyramids\_id\_2\_not\_every} fixes $q = 2$, as the
clause does. $h > 0$: same, but \texttt{grid\_pyramids\_id\_2\_knot} and
\texttt{grid\_pyramids\_id\_2\_coupling} do not require it: weaker. $w$ a mask: same,
\texttt{IsGridMask}. (2): for the knot kernel, $w$ infinitely divisible with the moment and a
family carrying the ladder: weaker, \texttt{grid\_pyramids\_id\_2\_knot} requires the carrying family
and a folded \Levy{} measure $\varpi_w$ of $w$, not the moment; the folded data of $\mu_{0,1}$:
same, \texttt{IsFoldedData}; for the carrier, \texttt{grid\_pyramids\_id\_2\_carrier} requires
$\varpi_w$, not the moment: weaker; without the moment, $w$ infinitely divisible: weaker,
\texttt{grid\_pyramids\_id\_2\_no\_kernel} requires only $\varpi_w$, which makes $w$ infinitely
divisible; $\int\log_+v\;\varpi_w(dv) = \infty$: same; for the uniform coupling, $w$ infinitely
divisible with the moment: same, as $\varpi_w$ with the moment. (3): the mask: same in
\texttt{grid\_pyramids\_id\_3}; $\varrho \ne \delta_0$ carried by $\xi\ZZ$, $\xi > 0$: same;
$\varrho$ symmetric and a probability law: weaker, \texttt{grid\_pyramids\_id\_3\_general} requires
neither; a fully covariant family: same, a \texttt{CascadeFamily} with its kernels. (4):
$0 < a < 1$: same; bounded support: same, $w\{|x| > R\} = 0$ for some $R$, in
\texttt{grid\_pyramids\_id\_4\_bounded}.)}
\end{proposition}

% CHANGED (wave 3e, 2026-11-15): the proof follows the machine-checked route. The line is read in
% R^1 through x -> x e_1 wherever a node of the next chapter is used, and back through y -> y_1;
% (1) reads the folded data of the mask from lem:line-exponent-data(1) and the lattice lemma on its
% lift; (2) identifies the knot data through the symmetric lift (a folded measure is determined by
% its lift), reads the family of prop:ladder-knot-kernel(2) back on the line, and computes the
% uniform coupling's kernels from (I1); the frequency clause bounds log_+ v by log_+(hv) +
% log_+(1/h); (4) takes the logarithm series through Re log(1 - a e^{i theta}) and checks that the
% sampled Laplace mask is a mask. Statement unchanged.
% CHANGED (fourth blind pass, 2026-11-15): the uniform coupling of (2) is built from the statement
% of lem:synthesis in R^1 (now in the uses), with the Levy measure of the knot data, the gluing and
% the folded data of its kernels proved here, where the proof had read it off the proof of
% prop:ladder-knot-kernel(2); the knot data are identified by lem:line-exponent-data(1) on the line
% in place of a uniqueness lem:family-structure does not state; (1) names the triplet that
% lem:lattice-infinitely-divisible reads; lem:family-dilations, cited for nothing, leaves the proof
% and its uses; (4) says that every mask of bounded support has no exact grid pyramid.
% CHANGED (length take (b), 2026-11-16): wording tightened, route and steps unchanged.
% CHANGED (the theory paper, 2026-11-16): the closing sentence on checking and cited facts shortened; no step changed.
% CHANGED (the theory paper's blueprint round, 2026-11-16, review of Section 2, PROSE): the
% identification of R with R^1 is stated once, in the lead of the section, and the proof no longer
% carries each law to R^1 and back; the symmetric lift of a folded measure, which has content, stays
% as its symmetrization varpi_s. No step changed.
\begin{proof}
\leanok
\uses{lem:line-axioms-kernels, lem:lattice-infinitely-divisible, prop:no-bounded-step,
prop:ladder-knot-kernel, lem:synthesis, lem:discrete-covariance-line, thm:lattice-line,
prop:no-lattice-kernel-line, prop:recursion-helmholtz,
prop:ladder-not-every-frequency, lem:line-exponent-data, lem:family-structure}
A measure $\varpi$ on $(0,\infty)$ has the symmetrization $\varpi_s$, the symmetric measure on
$\RR$ that puts half of $\varpi$ on each side of $0$. Then $\varpi_s(B) = \tfrac12\varpi(B)$ for
$B \subseteq (0,\infty)$, so a measure on $(0,\infty)$ is determined by its symmetrization, and
symmetrization commutes with sums and with dilations. Since $\cos$ is even,
$\int(1 - \cos\omega x)\,\varpi_s(dx) = \int_{(0,\infty)}(1 - \cos\omega v)\,\varpi(dv)$. So
$(a, \varpi)$ is a folded pair of the exponent of a law $m$ on $\RR$ if and only if $m$ has the
data $(a, \varpi_s)$.

(1) Let $w$ be infinitely divisible, with folded data $(a, \varpi_w)$
(Lemma~\ref{lem:line-exponent-data}(1)). Then $w$, carried by the lattice $h\ZZ$, has the data
$(a, (\varpi_w)_s)$, the generating triplet $(2a, (\varpi_w)_s, 0)$ in Sato's normalization
(Lemma~\ref{lem:symmetric-levy-khintchine}). By
Lemma~\ref{lem:lattice-infinitely-divisible} at $d = 1$, $2a = 0$, and $(\varpi_w)_s$ is finite and
carried by $h\ZZ \setminus \{0\}$. Hence $\varpi_w = 2(\varpi_w)_s$ on $(0,\infty)$ is finite and
carried by $h\ZZ_{>0}$. So $-\log\widehat w(\omega) = \int(1 - \cos\omega v)\,\varpi_w(dv)$, the
display, and $\widehat w > 0$. And $\varpi_w \ne 0$, since otherwise $w = \delta_0$.
Conversely, let $\widehat w$ have this form with $\varpi_w$ finite. Put
$|\varpi_w| := \varpi_w\bigl((0,\infty)\bigr)$ and let $\sigma$ be the symmetric law
$|\varpi_w|^{-1}(\varpi_w)_s$ on $\RR$. Then $\widehat w = e^{|\varpi_w|(\widehat\sigma - 1)}$
is the transform of the compound Poisson law
$\sum_{k \ge 0}e^{-|\varpi_w|}\tfrac{|\varpi_w|^k}{k!}\sigma^{*k}$, and the $n$-th power of
$e^{(|\varpi_w|/n)(\widehat\sigma - 1)}$, a transform of the same kind, for every $n$. So $w$ is
infinitely divisible. The mask, symmetric and other than $\delta_0$, is not a point mass. So by
Proposition~\ref{prop:no-bounded-step} at $d = 1$ an infinitely divisible mask has unbounded
support, and a mask of bounded support is not infinitely divisible.

(2) Suppose the grid pyramid is exact, carried by a $q^{\ZZ}$-covariant family in the
canonical discrete gauge. By Lemma~\ref{lem:line-axioms-kernels} and
Lemma~\ref{lem:discrete-covariance-line} its kernels are a symmetric family in dimension $1$,
$q$-covariant with $S_qt = qt$ and the step $\mu_{1,q} = w$. Its step $w$ is infinitely divisible
(Lemma~\ref{lem:family-structure}), so by (1) $w$ has the data $(0, (\varpi_w)_s)$.
Proposition~\ref{prop:ladder-knot-kernel}(1) at $d = 1$ gives
$\int\log_+|x|\,(\varpi_w)_s(dx) < \infty$, which is $\int\log_+v\;\varpi_w(dv) < \infty$. It
also gives the product formula, $A_1 = A_w/(q^2 - 1) = 0$ and
$\nu_1 = \nu^* = \sum_{m \ge 1}D_{q^{-m}}(\varpi_w)_s$, the symmetrization of
$\varpi_1 = \sum_{m \ge 1}T_{q^{-m}}\varpi_w$. By the first paragraph and the uniqueness of the
folded data (Lemma~\ref{lem:line-exponent-data}(1)), $a(1) = 0$ and $\mu_{0,1}$ has the folded
measure $\varpi_1$. These depend on $w$ only, so the knot kernel is unique. The partial products
are the transforms of the partial convolutions, by $\widehat{T_c m}(\omega) = \widehat m(c\omega)$
and the convolution theorem.

Conversely, let $w$ be infinitely divisible with the moment. By (1), $\varpi_w$ is finite and
nonzero, and $w$ is a symmetric infinitely divisible law on $\RR$, other than $\delta_0$, with the
data $(0, (\varpi_w)_s)$ and $\int\log_+|x|\,(\varpi_w)_s(dx) = \int\log_+v\;\varpi_w(dv) < \infty$.
Proposition~\ref{prop:ladder-knot-kernel}(2) at $d = 1$ gives a symmetric $q$-covariant family in
dimension $1$ with $S_qt = qt$ and the step $\mu_{1,q} = w$. By
Lemma~\ref{lem:line-axioms-kernels}(2) its laws are the kernels of a family with (A1)--(A7) and
(ND), which by Lemma~\ref{lem:discrete-covariance-line} satisfies (A8$|q^{\ZZ}$) in the canonical
discrete gauge. Covariance at $q$, iterated and inverted, gives
$\mu_{q^ns,q^nt} = T_{q^n}\mu_{s,t}$ for every $n \in \ZZ$, so
$\mu_{q^n,q^{n+1}} = T_{q^n}w$: the pyramid is exact. Proposition~\ref{prop:ladder-knot-kernel}(2)
at $d = 1$, applied to $w$, gives two families with the ladder of $w$ that differ at some pair
of scales. The same clause makes the moment of $(\varpi_w)_s$ equivalent to
$\int\log_+|x|\,w(dx) < \infty$.

The measure $\varpi_w \ne 0$ is carried by the countable set $h\ZZ_{>0}$, so it has an atom
$v_0 \in h\ZZ_{>0}$, and $T_{q^{-m}}\varpi_w$ is carried by $q^{-m}h\ZZ_{>0}$. So $\varpi_1$ is
carried by the union of these sets, with the atoms $q^{-m}v_0 \to 0$. A measure carried by $\xi\ZZ$
has no atom in $(0, \xi)$, so $\varpi_1$ is carried by no lattice.

Let $w$ be infinitely divisible without the moment. Then
Proposition~\ref{prop:ladder-knot-kernel}(3) at $d = 1$, applied to $w$, says that the partial
products have no limit that is a characteristic function, and that they tend to $0$ at almost
every $\omega$. No family
carries the ladder, as one would force the moment by the ``only if'' direction above. The partial
products need not tend to $0$ at every $\omega \ne 0$. Take $q = 2$, let $\varrho$ be the law on
$\ZZ$ of Proposition~\ref{prop:ladder-not-every-frequency}, and put
$w := T_h\varrho$. Its transform $\widehat w(\omega) = \widehat\varrho(h\omega)$ has the form of
(1) with the finite measure $\varpi_w = \sum_{j \ge 1}j^{-2}\delta_{h2^j} \ne 0$, so $w$ is
infinitely divisible, symmetric, other than $\delta_0$ and carried by $h\ZZ$. The moment fails:
$\log_+v \le \log_+(hv) + \log_+(1/h)$, so
$\int\log_+(hv)\,\varpi_\varrho(dv) \ge \int\log_+v\;\varpi_\varrho(dv) -
\log_+(1/h)\,\varpi_\varrho\bigl((0,\infty)\bigr) = \infty$, with $\varpi_\varrho = \sum_{j \ge
1}j^{-2}\delta_{2^j}$ of infinite moment by that proposition. The partial products of $w$ at
$\omega = 2\pi/h$ are those of $\varrho$ at $2\pi$, which have a positive limit by that
proposition.

(3) By Lemma~\ref{lem:line-axioms-kernels} the kernels of a fully covariant family form a
symmetric family in dimension $1$ that is $\lambda$-covariant for every $\lambda > 0$, and $\xi\ZZ$
is a full-rank lattice of $\RR$. Proposition~\ref{prop:ladder-knot-kernel}(4)(a) at $d = 1$ is
then the second sentence. It implies
the first, as a grid pyramid exact in the third stratum would have the step $\mu_{1,q} = w$, a law
other than $\delta_0$ carried by $h\ZZ$. A family carrying an exact grid pyramid is in the second
stratum by definition, and by the second sentence it is not fully covariant, its step $w$ being a
lattice law.

(4) The point masses $\tfrac{1-a}{1+a}a^{|n|}$ sum to $1$, since $\sum_na^{|n|} = \tfrac{1+a}{1-a}$.
The transform of $\lambda^{\mathrm{L}}_a$ is the geometric series
$\tfrac{1-a}{1+a}\sum_na^{|n|}e^{inh\omega} = (1-a)^2/(1 - 2a\cos h\omega + a^2)$
(Proposition~\ref{prop:recursion-helmholtz}(2), at the spacing $1$ and the range $t$ with
$e^{-1/t} = a$). Put $z := ae^{i\vartheta}$. Then $1 - 2a\cos\vartheta + a^2 = |1 - z|^2$, and the
series of the logarithm gives
$-\log|1 - z| = \operatorname{Re}\sum_{k \ge 1}z^k/k = \sum_{k \ge 1}\tfrac{a^k}{k}\cos k\vartheta$,
which at $\vartheta = 0$ is $-\log(1 - a)$. Hence
$\log\tfrac{1 - 2a\cos\vartheta + a^2}{(1-a)^2} = \sum_{k \ge 1}\tfrac{2a^k}{k}(1 - \cos k\vartheta)$,
and at $\vartheta = h\omega$
\[
  \widehat{\lambda^{\mathrm{L}}_a}(\omega) = \frac{1-a}{1 - ae^{ih\omega}}\cdot\frac{1-a}{1 - ae^{-ih\omega}}
  = \exp\Bigl(-\sum_{k \ge 1}\frac{2a^k}{k}\bigl(1 - \cos(kh\omega)\bigr)\Bigr) .
\]
So $\lambda^{\mathrm{L}}_a$ has the form of (1) with the finite measure
$\varpi = \sum_{k \ge 1}\tfrac{2a^k}{k}\delta_{kh} \ne 0$ on $h\ZZ_{>0}$; it is therefore symmetric,
other than $\delta_0$, a mask, and infinitely divisible. Its moments are finite for every $p > 0$:
$|hn|^p\tfrac{1-a}{1+a}a^{|n|} \le h^p|n|^Ka^{|n|}$ with $K$ the least integer $\ge p$, and
$\sum_n|n|^Ka^{|n|} < \infty$. As $\log_+x \le x$ for $x \ge 0$, the logarithmic moment is finite,
and (2) applies. As $1 - 2a\cos\vartheta + a^2 = (1-a)^2 + 2a(1 - \cos\vartheta)$, the transform is
$[1 + 2\varsigma(1 - \cos h\omega)]^{-1}$, and the product formula of (2) is the display. Clause
(3) gives the stratum. The binomial mask has the transform $\cos^4(h\omega/2)$, by the binomial
theorem: it is even, so the mask is symmetric, and it vanishes at $\omega = \pi/h$, so the mask is
not $\delta_0$; its support is $\{0, \pm h, \pm2h\}$. A mask of bounded support is not infinitely
divisible by (1), so its grid pyramid is not exact by (2). The binomial masks and the sampled B-spline masks
are masks of bounded support. The binomial weights $4^{-k}\binom{2k}{k+n}$ are symmetric in
$n$, sum to $1$ and are not concentrated at $0$ for $k \ge 1$. A B-spline of degree $r \ge 2$
with unit knot spacing, centred at $0$, is even and positive exactly on $(-(r+1)/2, (r+1)/2)$, so its
values at the integers, placed at $hn$ and renormalized, form a symmetric law with finitely many
atoms, among them $0$ and $\pm h$.

% CHANGED (fidelity review, 2026-11-15, R64): the closing sentence names the one part of (4) that is
% not machine-checked; it said "Every step is machine-checked".
Every step is machine-checked but one: in (4), that the binomial masks other than
$\tfrac1{16}[1,4,6,4,1]$ and the B-spline samples are symmetric masks of bounded support is shown on
paper, in the paragraph above. That a mask of bounded support has no exact grid pyramid is checked
for every such mask.
\end{proof}

The covariant kernel of a grid pyramid is the equivalent weighting function of Burt and Adelson
[\sscite{burt1983laplacian}, p.~533], run downward in scale: $T_{q^{-1}}w * \dots * T_{q^{-M}}w$ is
the equivalent weight of $M$ levels, rescaled back to the scale of the step $w$, and its limit is a continuum
kernel, as Proposition~\ref{prop:no-lattice-kernel-line} requires. The families that carry the
ladder are the couplings of the ladder in the sense of Appendix~\ref{ch:line-appendix}, how each
atom of the mask is switched on across one octave. The mask $\lambda^{\mathrm{L}}_a$ is compound
Poisson with the rate $\varpi\bigl((0,\infty)\bigr) = \sum_k2a^k/k = 2\log\tfrac{1}{1-a}$. At
$a = e^{-h/t}$, the sampled Laplace kernel of range $t$, the rate is
$2\log(t/h) + O(h/t)$. So the expected number of lattice jumps that make up one step grows without
bound as the grid refines relative to the scale. The continuum \Matern{} pyramid
(Theorem~\ref{thm:matern-pyramid-line}) has one Laplace kernel per octave, and on the grid that one
kernel is a compound Poisson law with this many jumps on average.

Proposition~\ref{prop:grid-pyramids-infinitely-divisible}(3) says that the mask is not a step of a
fully covariant family. It does not say that the arrays of a grid pyramid are not the samples of
one. For the sampled Laplace mask the computed level is the sampled \Matern{} level at the first
level only, although every step is the sampled Laplace step.
Proposition~\ref{prop:sampled-laplace-exact}(2) samples the \Matern{} pyramid on the grids
$q^nh\ZZ$ and finds the sampled Laplace step at every level. From the second level on, the array
the steps compute is the restriction of a convolution of sampled Laplace kernels, not the sample of
their convolution, and it is not the sampled \Matern{} level. The two readings of the same arrays,
steps of a family in the second stratum and samples of one in the third, agree at level $1$ and
differ from level $2$. The computation is the first: the level-$n$ array is the ladder step
$w * T_qw * \dots * T_{q^{n-1}}w$ of the geometric pyramid read on the grid $q^nh\ZZ$
(Proposition~\ref{prop:line-and-arrays}(2)).

\subsection{Grids on the line: the sampled Laplace family}
\label{subsec:sampled-laplace}

On the line the \Matern{} family is computed by forward and backward first-order sections
\prog{V-cone v0.1}{Prop.~7.2, pp.~60--61} (Theorem~\ref{thm:cited-A19}). This subsection holds what is specific to grids. The sampled Laplace family is
an exact cascade on the grid. The \Matern{} pyramid computed on the grids agrees with its sampled
steps at the first level, and from the second level the computed arrays are not the sampled levels. That the level
differences of the \Matern{} pyramid are Laplacians is Proposition~\ref{prop:laplacian-pyramid}, in
every dimension.

% The split (adr/0007, 2026-11-16): moved here verbatim from the computation chapter (its
% section 6.1), since prop:sampled-laplace-exact, prop:discrete-helmholtz and
% thm:matern-pyramid-line read it, and its Lean file is imported by the theory's. The computation
% paper cites it behind.
% draft: Proposition 6.1 (section 6.1), lines 1057-1074 of draft/spatial-hemigroup-pyramids.md.
% Moved from the affine blueprint (06-paper-vi.tex at origin/main, prop:recursion-helmholtz), plan
% decision 5; statement and proof are the affine node's, which the draft prints unchanged.
% CHANGED: renamed the Laplace density ell_t to L_t (BLUEPRINT-INPUTS: ell is log q, defined once in
% chapter 2), in the statement and the proof.
% CHANGED: in the proof, renamed c := max(0, x) to x_+ and k := a/(1-a)^2 to varsigma
% (BLUEPRINT-INPUTS).
% CHANGED (blind review 2026-11-12, flag 24): in the proof, the integration variables y and u are
% written xi, u = h/2t is beta, and the auxiliary sequence of the factorization check is z, so that
% y is only the input, u only a sequence and w only the forward recursion. Some semicolons of the
% proof become full stops (flag 34). No step changed.
\begin{proposition}[the recursion is a Helmholtz solve]
\label{prop:recursion-helmholtz}
\lean{Pyramids.recursion_line_transforms, Pyramids.recursion_line_passes,
Pyramids.recursion_line_increment, Pyramids.recursion_grid_transform,
Pyramids.recursion_grid_green, Pyramids.recursion_grid_factor,
Pyramids.recursion_grid_passes}\leanok
\begin{enumerate}
\item \emph{(The line.)} For $t > 0$ let $L_t(x) := (2t)^{-1}e^{-|x|/t}$, the Laplace density,
and $e^+_t(x) := t^{-1}e^{-x/t}\mathbf 1_{\{x > 0\}}$, $e^-_t(x) := e^+_t(-x)$. With
$\widehat f(\omega) = \int e^{i\omega x}f(x)\,dx$,
\[
  \widehat{e^\pm_t}(\omega) = (1 \mp it\omega)^{-1}, \qquad
  \widehat{L_t}(\omega) = (1 + t^2\omega^2)^{-1}, \qquad L_t = e^+_t * e^-_t .
\]
For $y \in L^1(\RR)$, $w := e^+_t * y$ is locally absolutely continuous with
$(1 + t\partial_x)w = y$ almost everywhere, and $e^-_t * y$ likewise solves
$(1 - t\partial_x)v = y$: convolution with $L_t$ is a forward pass followed by a backward pass.
% CHANGED (2026-11-14, a strengthening): gamma >= 1 widened to gamma >= 0; the proof never read it.
For $0 < s < t$ and an integer $\gamma \ge 0$,
\[
  \Bigl(\frac{1 + s^2\omega^2}{1 + t^2\omega^2}\Bigr)^{\gamma}
   = \Bigl(\frac{s^2}{t^2} + \Bigl(1 - \frac{s^2}{t^2}\Bigr)\widehat{L_t}(\omega)\Bigr)^{\gamma} ,
\]
the transform of the $\gamma$-fold convolution power of the probability measure
$\tfrac{s^2}{t^2}\delta_0 + (1 - \tfrac{s^2}{t^2})L_t(x)dx$: the \Matern{} increment of index
$\gamma$ from $s$ to $t$ is the $\gamma$-th convolution power of one convex combination of the
identity with one forward--backward pair.
\item \emph{(The grid.)} Fix $h > 0$, put $a := e^{-h/t}$ and let
$\lambda_t(n) := \frac{1-a}{1+a}\,a^{|n|}$, $n \in \ZZ$, the Laplace kernel sampled on $h\ZZ$ and
renormalized to unit mass. With $\widehat\lambda(\vartheta) = \sum_n\lambda(n)e^{in\vartheta}$,
$D_h(\vartheta) := h^{-2}(2 - 2\cos\vartheta)$ and $\Delta_hu(n) := h^{-2}(u(n+1) - 2u(n) + u(n-1))$,
\[
  \widehat{\lambda_t}(\vartheta) = \frac{(1-a)^2}{1 - 2a\cos\vartheta + a^2}
   = \frac{1}{1 + t_h^2D_h(\vartheta)} ,
  \qquad t_h := \frac{h}{2\sinh(h/2t)} ,
\]
so $\lambda_t$ is the unique summable solution of $(1 - t_h^2\Delta_h)\lambda_t = \delta_0$, the
discrete Green's function of the three-point Helmholtz operator at the renormalized range $t_h$;
$t_h < t$ and $t_h/t \to 1$ as $h/t \to 0$. With $(Su)(n) := u(n-1)$,
\[
  1 - t_h^2\Delta_h = (1-a)^{-2}\,(1 - aS)(1 - aS^{-1})
\]
as operators on $\ell^1(\ZZ)$, and $\lambda_t * y = (1-a)^2(1-aS^{-1})^{-1}(1-aS)^{-1}y$ is the
forward recursion $w(n) = y(n) + a\,w(n-1)$ followed by the backward recursion
$v(n) = w(n) + a\,v(n+1)$, scaled by $(1-a)^2$: the lower and upper bidiagonal factors of the
tridiagonal operator.
\end{enumerate}
% CHANGED (fidelity review, 2026-11-15, R26/R50/R69): the status line lists the declarations' hypotheses beside the node's (LINKAGE rule 6).
\statusT\quad\emph{(Cites nothing beyond the Lebesgue differentiation theorem on the line.
Machine-checked, Lean core, no ledger entry. \textbf{The declarations.}
\texttt{recursion\_line\_transforms} carries the three transforms and $L_t = e^+_t * e^-_t$ at every
point. Hypothesis: $t > 0$, same. \texttt{recursion\_line\_passes} carries the two passes: each
is absolutely continuous on every bounded interval and solves its equation almost everywhere, with
the derivative read pointwise. Hypotheses: $t > 0$ and $y$ real-valued and integrable, same.
\texttt{recursion\_line\_increment} carries that the mixture is a probability measure, the identity
of transforms, and the transform of the $\gamma$-fold convolution power. Hypotheses: $0 < s < t$,
% CHANGED (2026-11-14, a strengthening): the unused gamma >= 1 dropped.
and any integer $\gamma \ge 0$, same (the unused $\gamma \ge 1$ was dropped on 2026-11-14).
\texttt{recursion\_grid\_transform} carries both forms of
$\widehat{\lambda_t}$. Hypotheses: $h, t > 0$, same. \texttt{recursion\_grid\_green} carries
that $\lambda_t$ is summable, solves the equation at every $n$ and is the only summable solution,
that $t_h < t$, and that $t_h/t \to 1$ as $h \to 0$ from above at fixed $t$. That is the
statement's limit, since $t_h/t$ is a function of $h/t$. Hypotheses: $h, t > 0$, same.
\texttt{recursion\_grid\_factor} carries the factorization on every sequence on $\ZZ$, stronger
than on $\ell^1(\ZZ)$. Hypotheses: $h, t > 0$, same. \texttt{recursion\_grid\_passes} carries
that the two recursions have unique summable solutions $w$ and $v$ and that
$\lambda_t * y = (1-a)^2v$ at every $n$. Hypotheses: $h, t > 0$ and $y$ summable, same.
\textbf{Not claimed}: the reading of $\lambda_t$ as the periodization of the continuous transfer.
\textbf{Hypotheses} (the node's, each with what its declarations require): (1): $t > 0$: same;
$y \in L^1(\RR)$: same for real-valued $y$, which \texttt{recursion\_line\_passes} requires
(\texttt{Integrable} on $\RR \to \RR$); a complex $y$ is its real and imaginary parts; $0 < s < t$:
same; $\gamma \ge 0$ an integer: same. (2): $h > 0$ and $t > 0$: same, in each of the four grid
declarations; $y \in \ell^1(\ZZ)$: same, real-valued (\texttt{Summable} on $\ZZ \to \RR$); the
factorization on $\ell^1(\ZZ)$: weaker, \texttt{recursion\_grid\_factor} requires nothing of the
sequence.)}
\end{proposition}

\begin{proof}
\leanok
\emph{Clause (1).} $\int_0^\infty t^{-1}e^{-x/t}e^{i\omega x}dx = t^{-1}\int_0^\infty
e^{(i\omega - 1/t)x}dx = (1 - it\omega)^{-1}$, and $e^-_t$ is the reflection of $e^+_t$, so its
transform is $(1 + it\omega)^{-1}$. Off the null set $\{0\}$, $L_t = \tfrac12(e^+_t + e^-_t)$, so
$\widehat{L_t}(\omega) = \tfrac12\bigl((1 - it\omega)^{-1} + (1 + it\omega)^{-1}\bigr) =
(1 + t^2\omega^2)^{-1}$. For every $x$, with $x_+ := \max(0, x)$,
\[
  (e^+_t * e^-_t)(x) = \int_{x_+}^\infty t^{-2}e^{-\xi/t}e^{-(\xi-x)/t}\,d\xi
   = t^{-2}e^{x/t}\,\tfrac t2\,e^{-2x_+/t} = (2t)^{-1}e^{(x - 2x_+)/t} = L_t(x) ,
\]
since $x - 2x_+ = -|x|$. For $y \in L^1$, writing out
$e^+_t(x - z) = t^{-1}e^{-(x-z)/t}\mathbf 1_{\{z < x\}}$ gives
$w(x) = (e^+_t * y)(x) = t^{-1}e^{-x/t}F(x)$ with $F(x) := \int_{-\infty}^xe^{\xi/t}y(\xi)\,d\xi$, where
$\xi \mapsto e^{\xi/t}y(\xi)$ is integrable on every $(-\infty, b]$. So
$F(x) = F(b) + \int_b^xe^{\xi/t}y(\xi)\,d\xi$ is absolutely continuous on every interval, and so is
$w$, a product with a smooth function. By the Lebesgue differentiation theorem
$F'(x) = e^{x/t}y(x)$ for almost every $x$, and there $w'(x) = -t^{-1}w(x) + t^{-1}y(x)$, that is
$w + tw' = y$. Likewise $(e^-_t * y)(x) = t^{-1}e^{x/t}\int_x^\infty e^{-\xi/t}y(\xi)\,d\xi$, which
gives $v - tv' = y$ almost everywhere. The increment identity is
$1 + s^2\omega^2 = \tfrac{s^2}{t^2}(1 + t^2\omega^2) + (1 - \tfrac{s^2}{t^2})$ divided by
$1 + t^2\omega^2$ and raised to the power $\gamma$. The measure is positive with mass $1$ because
$s < t$ and $\int L_t = \widehat{L_t}(0) = 1$, its transform is
$\tfrac{s^2}{t^2} + (1 - \tfrac{s^2}{t^2})\widehat{L_t}$, and the transform of a convolution power
is the power of the transform.

\emph{Clause (2).} Summing the two geometric series over $n \ge 0$ and $n < 0$,
$\sum_na^{|n|}e^{in\vartheta} = (1 - ae^{i\vartheta})^{-1} + ae^{-i\vartheta}(1 - ae^{-i\vartheta})^{-1}
= (1-a^2)/(1 - 2a\cos\vartheta + a^2)$, which gives the first form. Next
$1 - 2a\cos\vartheta + a^2 = (1-a)^2 + a(2 - 2\cos\vartheta)$ and
$t_h^2/h^2 = (2\sinh(h/2t))^{-2} = (a^{-1/2} - a^{1/2})^{-2} = a/(1-a)^2$, which gives the second.
With $\varsigma := t_h^2/h^2 = a/(1-a)^2$,
$(1 - t_h^2\Delta_h)u(n) = u(n) - \varsigma\,(u(n+1) - 2u(n) + u(n-1))$.
For $\lambda_t(n) = c\,a^{|n|}$, $c := (1-a)/(1+a)$, this is
$c\,a^{n-1}(a - \varsigma(1-a)^2) = 0$ for $n \ge 1$, the same for $n \le -1$, and
$c\,(1 + 2\varsigma(1-a)) = c\,(1+a)/(1-a) = 1$ at $n = 0$, and $\lambda_t$ is summable. A bounded
solution $e$ of the homogeneous equation satisfies $e(n+1) + e(n-1) = (a + a^{-1})e(n)$, so
$f(n) := e(n+1) - a\,e(n)$ has $f(n) = a^{-1}f(n-1)$. Then $f(m) = a^{-m}f(0)$ is bounded for
$m \ge 0$ only if $f(0) = 0$, and then $f = 0$, that is $e(n+1) = a\,e(n)$. Now $e(0) = a^me(-m)$ with
$e(-m)$ bounded forces $e(0) = 0$, and then $e = 0$. The difference of two summable solutions is
bounded, so $\lambda_t$ is the unique summable solution. $\sinh x > x$ for $x > 0$ gives $t_h < t$,
and $t_h/t = \beta/\sinh\beta$ with $\beta := h/2t$, which tends to $1$ as $\beta \to 0$ because the
derivative of $\sinh$ at $0$ is $\cosh 0 = 1$. The factorization is checked on every sequence $u$:
with $z(m) := u(m) - a\,u(m+1)$, $z(n) - a\,z(n-1) = (1 + a^2)u(n) - a(u(n+1) + u(n-1))$, and
$(1-a)^{-2}$ times this is $u(n) - \varsigma(u(n+1) - 2u(n) + u(n-1))$. For $y \in \ell^1(\ZZ)$ and
$0 < a < 1$, $w(n) := \sum_{j \ge 0}a^jy(n - j)$ is summable, since
$\sum_n\sum_ja^j|y(n-j)| = (1-a)^{-1}\sum_n|y(n)|$, and solves $w(n) = y(n) + a\,w(n-1)$. The
difference $e$ of two summable solutions has $e(n) = a^me(n - m)$, so $a^{-m}e(n)$ stays bounded and
$e(n) = 0$. The backward recursion $v(n) := \sum_{j \ge 0}a^jw(n+j)$ is the same with the order of
$\ZZ$ reversed. Finally $\lambda_t * y$ and $(1-a)^2v$ are both bounded solutions of
$(1 - t_h^2\Delta_h)u = y$: the first because $(1 - t_h^2\Delta_h)(\lambda_t * y) =
((1 - t_h^2\Delta_h)\lambda_t) * y = \delta_0 * y$, the shifts passing through the convolution, and
$|\lambda_t * y| \le \sum|y|$. The second is one by the factorization, since
$(1-a)^2v(m) - a(1-a)^2v(m+1) = (1-a)^2w(m)$ and $w(n) - a\,w(n-1) = y(n)$. Their difference is a
bounded solution of the homogeneous equation, hence $0$.
\end{proof}

% draft: Proposition 2.23 (section 2.7), lines 303-313.
% CHANGED: the identity with the periodization is stated with its constant C_t and proved by
% Fourier series; the draft says "up to the normalization" and "verified numerically".
% CHANGED: in (1) the brackets of (2.13) are renormalized, p in place of a_s/a_t, so that each is
% the convex combination the draft's sentence names; the identity is the same. The draft's
% <check> on the gamma-fold powers is editorial and is dropped; the clause claims nothing about
% samples.
% CHANGED: in (1) the sentence "The sampled Gaussian family has no such property: its increments
% are not positive grid kernels, which is why the discrete Gaussian of [Lindeberg 1990] had to
% replace it" leaves the statement. It rested on two computed scale pairs (the blind review's
% flag, records/pyramids/PLAN-review-response.md, answered by qualifying it to the fine scales and
% moving the numbers to a labelled measurement); it is the measured paragraph after the proof.
% CHANGED: (2) defines the computation and the sampled level, and states the failure from the
% second level for every n >= 2, which the proof proves. The draft's "the error is roughly
% constant across levels, covariance holding h/(ct) fixed, and falls as h/(ct) does" is a
% measurement and leaves the statement for the measured paragraph (narrowing).
% CHANGED: (3) is stated as the identity with the normalization and the bound 1/12, which is
% what the draft's "relative size O((h/t)^2) at omega = 0" means; the draft's <check> is closed.
% CHANGED (blind review, flags 3, 4, 6, 19): the computed array is U_n and the sampled level V_n
% (S was also the bijection S_lambda, C also C_t); the recursion output is v; "an exact cascade
% family" is "an exact hemigroup on the grid", glossed in the section lead. The draft's (2.12) and
% (2.13) carry the labels eq:sampled-laplace-transfer and eq:sampled-laplace-increment.
% CHANGED (second blind review, L3): a short title; in (2) the ratio e^{-h/(ct)} is written
% a_{ct}, in the notation a_t of the statement, and the letter a alone is used only for it in the
% proof; V_n is "the sampled step", the samples of the step from t to q^n t. Renames.
% Moved from the status (second blind review, L7): the draft's proof of (2) for n >= 2 was the
% sentence "which differ" with a numerical table; the comparison of coefficients replaces it, and
% the table is the measured paragraph after the proof.
\begin{proposition}[the sampled Laplace family on the grid]
\label{prop:sampled-laplace-exact}
\lean{Pyramids.sampled_laplace_exact_sample, Pyramids.sampled_laplace_exact_transfer,
Pyramids.sampled_laplace_exact_periodization, Pyramids.sampled_laplace_exact_1_factor,
Pyramids.sampled_laplace_exact_1_hemigroup, Pyramids.sampled_laplace_exact_1_recursion,
Pyramids.sampled_laplace_exact_1_powers, Pyramids.subsampling_sampled_laplace,
Pyramids.sampled_laplace_exact_2_array, Pyramids.sampled_laplace_exact_2_restriction,
Pyramids.sampled_laplace_exact_2_first, Pyramids.sampled_laplace_exact_2_later,
Pyramids.sampled_laplace_exact_3}
\leanok
% CHANGED (the split, adr/0007, 2026-11-16): prop:helmholtz-subsampling leaves the uses. The
% subsampling sentence of (2) is this node's own declaration subsampling_sampled_laplace, whose
% checked proof is the identity beta^{q|m|} = (beta^q)^{|m|} that the proof below writes out; the
% citation of the computation paper's node was a parallel, not a step. Statement unchanged.
% CHANGED (the theory paper's blueprint round, 2026-11-16, review of Section 2, NOTATION): each
% clause names its sense of def:exact-senses (now in the uses): (1) "an exact hemigroup on the
% grid" reads "an exact cascade on the grid", the name the prose after that definition gives a
% family on one grid; in (2) the subsampling sentence is named sense (ii), and "the first level is
% exact" reads "the computation is exact in sense (i) at the first level, and at no later one".
% The titles follow. Names only; no clause changed.
\uses{prop:recursion-helmholtz, thm:matern-pyramid-line, def:exact-senses}
Fix $h > 0$. For $t > 0$ put $a_t := e^{-h/t}$, put $a_0 := 0$, and let
$\lambda_t(n) := \tfrac{1-a_t}{1+a_t}a_t^{|n|}$, $n \in \ZZ$: the Laplace kernel
$L_t(x) = (2t)^{-1}e^{-|x|/t}$ of range $t$, sampled on $h\ZZ$ and renormalized to unit mass, and
$\lambda_0 = \delta_0$. Its discrete transfer is
\begin{equation}
\label{eq:sampled-laplace-transfer}
  \widehat\lambda_t(\vartheta) := \sum_n\lambda_t(n)e^{in\vartheta}
  = \frac{(1-a_t)^2}{1 - 2a_t\cos\vartheta + a_t^2}
  = \frac{(1-a_t)^2}{(1 - a_te^{i\vartheta})(1 - a_te^{-i\vartheta})} ,
\end{equation}
and at $\vartheta = h\omega$ it is the periodization of the continuous transfer
$(1 + t^2\omega^2)^{-1}$ of $L_t$, up to a constant:
\[
  \widehat\lambda_t(h\omega) = \frac{1}{C_t}\sum_{n \in \ZZ}\frac{1}{1 + t^2(\omega + 2\pi n/h)^2},
  \qquad C_t := \frac{h}{2t}\cdot\frac{1 + a_t}{1 - a_t} .
\]
\begin{enumerate}
% CHANGED (fourth blind pass, 2026-11-15, a rename): the one-sided geometric sequences, written
% e^+_t and e^-_t (statement: "let e^+_t(n) := (1-a_t)a_t^n for n >= 0, e^+_t(n) := 0 for n < 0,
% and e^-_t(n) := e^+_t(-n)"), are \varepsilon^\pm_t, since e^\pm_t are the one-sided exponential
% densities of prop:recursion-helmholtz(1). Nothing else in the statement changes.
\item \emph{(An exact cascade.)} For $0 \le s < t$ put
$p := \tfrac{a_s(1-a_t)}{a_t(1-a_s)} \in [0,1)$, and let $\varepsilon^+_t(n) := (1-a_t)a_t^n$ for $n \ge 0$,
$\varepsilon^+_t(n) := 0$ for $n < 0$, and $\varepsilon^-_t(n) := \varepsilon^+_t(-n)$. Then
\begin{equation}
\label{eq:sampled-laplace-increment}
  \frac{\widehat\lambda_t}{\widehat\lambda_s}(\vartheta)
  = \prod_\pm\Bigl[p + (1-p)\frac{1 - a_t}{1 - a_te^{\pm i\vartheta}}\Bigr] ,
\end{equation}
the transfer of the probability
$\lambda_{s,t} := (p\delta_0 + (1-p)\varepsilon^+_t) * (p\delta_0 + (1-p)\varepsilon^-_t)$ on $\ZZ$. With
$\lambda_{t,t} := \delta_0$, the kernels $\lambda_{s,t}$ are positive, symmetric and of unit mass,
$\lambda_{0,t} = \lambda_t$, $\lambda_{r,t} = \lambda_{s,t} * \lambda_{r,s}$ for $r \le s \le t$,
$\widehat\lambda_{s,t}(\vartheta)$ is continuous in $(s,t)$, and $\lambda_{s,t} \ne \delta_0$ for
$s < t$: an exact cascade on the grid, and translation
covariant. Convolution with
$p\delta_0 + (1-p)\varepsilon^+_t$ maps a bounded sequence $y$ to the unique bounded solution $v$ of
\[
  v(n) = a_t\,v(n-1) + (1 - a_t + pa_t)\,y(n) - pa_t\,y(n-1),
\]
one forward recursion with one state, and the factor with $\varepsilon^-_t$ is the same recursion run
backward. For an integer $\gamma \ge 1$ the powers $\lambda^{*\gamma}_{s,t}$ form a cascade family
of the same kind, with $\gamma$ forward and $\gamma$ backward recursions per increment.
\item \emph{(The \Matern{} pyramid on the grids.)} Let $q \ge 2$ be an integer and $c, t > 0$.
The subsample $m \mapsto \lambda_T(qm)$ of the kernel of range $T$ on $h\ZZ$ is proportional to the
kernel of the same range on $q h\ZZ$: it is exact in sense (ii) of
Definition~\ref{def:exact-senses}. So in grid units the kernel of range $ctq^k$ on the grid
$q^kh\ZZ$ is the same array $g := \lambda_{ct}$, $g(m) = \tfrac{1-a}{1+a}a^{|m|}$ with
$a := a_{ct} = e^{-h/(ct)}$, at every $k \ge 0$. In the \Matern{} pyramid of
Theorem~\ref{thm:matern-pyramid-line} the step from $q^kt$ to $q^{k+1}t$ is the Laplace kernel
$L_{ctq^k}$, so the step from $t$ to $q^nt$ is $L_{ct} * L_{ctq} * \dots * L_{ctq^{n-1}}$. Let
$V_n$ be its samples on $q^nh\ZZ$, renormalized to unit mass, the \emph{sampled step}. The pyramid
is computed by
\[
  u_{k+1}(m) := (g * u_k)(qm), \qquad m \in \ZZ ,
\]
one forward and one backward recursion on the current grid $q^kh\ZZ$ followed by decimation, with
the same coefficients at every level. Let $U_n$ be $u_n$ computed from $u_0 = \delta_0$ and
renormalized to unit mass. Then $U_n$ is, up to a constant, the restriction to $q^nh\ZZ$ of the
convolution of the sampled factors, $L_{ctq^k}$ sampled on $q^kh\ZZ$ for $k < n$, while $V_n$ is the
sample of their convolution. The computation is exact in sense (i) at the first level,
$U_1 = V_1$, and at no later one: for every $n \ge 2$ and every $h, c, t > 0$, $U_n \ne V_n$.
% CHANGED (fourth blind pass, 2026-11-15, the author's decision): "for t > 0" added; at t = 0 the
% clause is undefined (alpha_0 diverges), and sampled_laplace_exact_3 takes 0 < t. It was: "With
% $\alpha_t(\omega) := ...$, the aliased mass," with no range for t.
\item \emph{(Aliasing.)} For $t > 0$, with
$\alpha_t(\omega) := \sum_{n \ne 0}\bigl(1 + t^2(\omega + 2\pi n/h)^2\bigr)^{-1}$, the aliased mass,
\[
  \widehat\lambda_t(h\omega) = \frac{(1 + t^2\omega^2)^{-1} + \alpha_t(\omega)}{1 + \alpha_t(0)} ,
  \qquad 0 < \alpha_t(0) < \frac{h^2}{12\,t^2} .
\]
\end{enumerate}
% CHANGED (wave 3f, 2026-11-15): every clause is machine-checked, so statement and proof are marked
% proved; the status line names the checked route of the periodization and of (2).
% CHANGED (2026-11-15, a strengthening): unused hypotheses dropped from the Lean of three
% declarations, listed at the end of the status line; the print keeps the node's setting.
% CHANGED (fidelity review, 2026-11-15, R26/R50/R69): the status line lists the declarations' hypotheses beside the node's (LINKAGE rule 6).
\statusT\quad\emph{(No ledger entry. The transfer and the forward--backward recursion of the
sampled kernel are Proposition~\ref{prop:recursion-helmholtz}(2), machine-checked; the subsampling
sentence of (2) is proved here, also machine-checked; the step
of the \Matern{} pyramid is Theorem~\ref{thm:matern-pyramid-line}(1), recalled to name $V_n$, which
is defined as the samples of the convolution itself. Everything else is proved here: the
periodization by Poisson's summation formula, (1) by the factorization and the uniqueness of a
summable kernel with a given transfer, (2) by writing both arrays as sums of geometric sequences
with the same ratios and comparing two coefficients, (3) by $\sum 1/n^2 = \pi^2/6$. Every clause is
machine-checked; Poisson's formula and $\sum 1/n^2 = \pi^2/6$ are theorems of the library, not
cited facts. \textbf{Not claimed}: the size
of $U_n - V_n$ (measured, not proved);
that the powers $\lambda^{*\gamma}_t$ for $\gamma \ge 2$ are samples of the continuum \Matern{}
kernels of index $\gamma$ (they are not, for the reason of (2)); anything about the sampled
Gaussian family, which is measured in the paragraph after the proof. \textbf{Hypotheses} (the
node's, each with what its declarations require; three declarations had unused hypotheses dropped on
2026-11-15): $h > 0$: same, but \texttt{sampled\_laplace\_exact\_2\_array} does not require it:
weaker. $t > 0$ for $\lambda_t$ as a sampled kernel, and $t \ge 0$ for the transfer: same in
\texttt{sampled\_laplace\_exact\_sample}; weaker in \texttt{sampled\_laplace\_exact\_transfer},
which holds for every real $t$, $a_t$ lying in $[0,1)$ for all $t$; $t > 0$ for the
periodization: same. (1): $0 \le s < t$, and $0 \le r \le s \le t$: same; $y$ bounded: same;
$\gamma \ge 1$ an integer: same. The transfers are sums over $\ZZ$ with no summability required: the
summability of the steps is a conclusion of \texttt{sampled\_laplace\_exact\_1\_hemigroup}, not a
hypothesis. (2): $q \ge 2$ an integer: same, but \texttt{sampled\_laplace\_exact\_2\_first} does not
require it: weaker; $c, t > 0$: same; $T > 0$ in the subsampling sentence: same, in
\texttt{subsampling\_sampled\_laplace}; $n \ge 2$: same, in
\texttt{sampled\_laplace\_exact\_2\_later}. (3): $t > 0$: same.)}
\end{proposition}

% CHANGED (wave 3f, 2026-11-15)
% The proof follows the machine-checked route. The periodization is Poisson's summation formula for
% f(v) = h L_t(hv) e^{i omega h v} in place of the Fourier coefficients of P and the uniform
% convergence of its series; the uniqueness of a summable kernel with a given transfer is proved by
% the coefficient formula; the restriction carries the constants c_k of the sampled factors; the
% computed level tracks the two last coefficients only and computes R_n directly; the sampled step
% is computed by L_r * L_s = (r^2 L_r - s^2 L_s)/(r^2 - s^2), three exponential integrals, in place
% of partial fractions and the uniqueness of transforms; the comparison uses the independence of
% geometric sequences and the monotonicity of sinh(u)/u in place of the Vandermonde matrix and two
% power series. Statement unchanged.
% CHANGED (fourth blind pass, 2026-11-15): the rename to \varepsilon^\pm_t; (2) says why the
% computation is one forward and one backward recursion with the same coefficients at every level;
% G_0 := \delta_0.
% CHANGED (length take (b), 2026-11-16): wording tightened, route and steps unchanged.
% CHANGED (the theory paper, 2026-11-16): the closing sentence on checking and cited facts shortened; no step changed.
\begin{proof}
\leanok
\uses{prop:recursion-helmholtz}
\emph{The transfer and the periodization.} The first form of $\widehat\lambda_t$ is
Proposition~\ref{prop:recursion-helmholtz}(2), the sum of two geometric series; at $t = 0$ both
sides are $1$. The second is
$1 - 2a_t\cos\vartheta + a_t^2 = (1 - a_te^{i\vartheta})(1 - a_te^{-i\vartheta})$. For the
periodization let $t > 0$ and $f(v) := h\,L_t(hv)\,e^{i\omega hv}$ for $v \in \RR$, a continuous
function. Its samples are $f(n) = \tfrac{h}{2t}a_t^{|n|}e^{inh\omega}$, so
$\sum_nf(n) = C_t\,\widehat\lambda_t(h\omega)$. Its Fourier transform is, by the substitution
$x = hv$ and the transform $(1 + t^2\omega^2)^{-1}$ of $L_t$
(Proposition~\ref{prop:recursion-helmholtz}(1)),
\[
  \int_{\RR}f(v)e^{-2\pi iv\xi}\,dv = \int_{\RR}L_t(x)e^{i(\omega - 2\pi\xi/h)x}\,dx
  = \frac{1}{1 + t^2(\omega - 2\pi\xi/h)^2} .
\]
Both decay like $|v|^{-2}$: $|f(v)| \le (t/h)v^{-2}$ for $|v| \ge 1$, as $e^{-y} \le 2/y^2$ for
$y > 0$, and the transform is at most $h^2/(\pi^2t^2\xi^2)$ once $|\xi| \ge h|\omega|/\pi + 1$, as
then $|\omega - 2\pi\xi/h| \ge \pi|\xi|/h$. The same bound at $-\omega$ makes the series
$P(\omega) := \sum_{n \in \ZZ}\bigl(1 + t^2(\omega + 2\pi n/h)^2\bigr)^{-1}$ converge. So Poisson's
summation formula gives $\sum_nf(n) = \sum_n\widehat f(n)$, which is $P(\omega)$ after
$n \mapsto -n$, and $C_t\,\widehat\lambda_t(h\omega) = P(\omega)$.

\emph{(3).} Splitting off the term $n = 0$, $P(\omega) = (1 + t^2\omega^2)^{-1} + \alpha_t(\omega)$,
and the aliased mass converges with $P$. As $\widehat\lambda_t(0) = 1$,
$C_t = P(0) = 1 + \alpha_t(0)$, and $\widehat\lambda_t(h\omega) = P(\omega)/P(0)$ is the display.
Each term of $\alpha_t(0)$ is positive and, for $n \ne 0$, below $h^2/(4\pi^2n^2t^2)$, so
$0 < \alpha_t(0) < \tfrac{h^2}{4\pi^2t^2}\sum_{n \ne 0}n^{-2} = \tfrac{h^2}{4\pi^2t^2}\cdot
\tfrac{\pi^2}{3} = \tfrac{h^2}{12t^2}$, by $\sum_{n \ge 1}n^{-2} = \pi^2/6$.

\emph{Kernels on $\ZZ$.} Summable sequences on $\ZZ$ form a commutative and associative convolution
algebra with unit $\delta_0$, the transfer of a convolution being the product of the transfers. A
summable sequence $\lambda$ is determined by its transfer: integrating
$\widehat\lambda(\vartheta)e^{-ik\vartheta}$ over $[-\pi,\pi]$ term by term gives $2\pi\lambda(k)$.

\emph{(1).} For $0 \le s < t$, $0 \le a_s < a_t < 1$, so $p \ge 0$, and $p < 1$ is
$a_s(1-a_t) < a_t(1-a_s)$, that is $a_s < a_t$; $p = 0$ when $s = 0$. For every $z$ with
$a_tz \ne 1$,
\[
  p + (1-p)\frac{1 - a_t}{1 - a_tz} = \frac{1 - a_t}{1 - a_s}\cdot\frac{1 - a_sz}{1 - a_tz} :
\]
multiplying by $(1 - a_s)(1 - a_tz)$, the left side becomes
$a_s(1-a_t)(1 - a_tz)/a_t + (a_t - a_s)(1-a_t)/a_t = (1-a_t)(1 - a_sz)$. The transfers of
$\varepsilon^\pm_t$ are $(1-a_t)\sum_{n \ge 0}a_t^ne^{\pm in\vartheta} = (1-a_t)/(1 - a_te^{\pm i\vartheta})$,
so the product of this identity at $z = e^{\pm i\vartheta}$ gives, by the second form of the
transfer, $\widehat\lambda_{s,t} = \widehat\lambda_t/\widehat\lambda_s$, also at $s = t$. The kernel
$\lambda_{s,t}$ is a convolution of two probabilities, as $p \in [0,1)$, and symmetric, the
convolution of a law with its reflection. By the uniqueness above, $\lambda_{0,t} = \lambda_t$
($\widehat\lambda_0 = 1$ as $a_0 = 0$) and $\lambda_{r,t} = \lambda_{s,t} * \lambda_{r,s}$, as
$\widehat\lambda_t/\widehat\lambda_r = (\widehat\lambda_t/\widehat\lambda_s)
(\widehat\lambda_s/\widehat\lambda_r)$. The map $t \mapsto a_t$ is continuous, being
$x \mapsto e^{-1/x}$ at $x = t/h$, extended by $0$ for $x \le 0$, and
$1 - 2a\cos\vartheta + a^2 \ge (1-a)^2 > 0$ for $0 \le a < 1$. So $\widehat\lambda_t(\vartheta)$
is positive and continuous in $t$, and the quotients are continuous in $(s,t)$. At
$\vartheta = \pi$ the transfer is $\bigl((1 - a_t)/(1 + a_t)\bigr)^2$, strictly decreasing in
$a_t$, so $\widehat\lambda_{s,t}(\pi) < 1$ and $\lambda_{s,t} \ne \delta_0$ for $s < t$.

For the recursion let $F := p\delta_0 + (1-p)\varepsilon^+_t$. Since
$\varepsilon^+_t(n) = (1 - a_t)\delta_0(n) + a_t\varepsilon^+_t(n-1)$ for every $n \in \ZZ$,
\[
  F(n) = (1 - a_t + pa_t)\,\delta_0(n) - pa_t\,\delta_0(n-1) + a_t\,F(n-1) .
\]
For a bounded $y$, $v := F * y$ converges absolutely and is bounded by $\sup|y|\sum|F|$, and
convolving the identity with $y$ gives the recursion. Two bounded solutions differ by a bounded
$e$ with $e(n) = a_t^ke(n-k)$ for every $k \ge 0$, so $e = 0$. The reflected factor $F(-\cdot)$
satisfies the identity with $n - 1$ replaced by $n + 1$, which gives the backward recursion in the
same way. For the powers, positivity, symmetry and unit mass pass to convolutions, hence by
induction to $\lambda^{*\gamma}_{s,t}$, whose transfer is $\widehat\lambda_{s,t}^\gamma$. By the
uniqueness, $(\lambda * \mu)^{*\gamma} = \lambda^{*\gamma} * \mu^{*\gamma}$ for summable $\lambda$
and $\mu$. This gives the composition rule for the powers and, with $\lambda_{s,t} = F * F(-\cdot)$,
the $\gamma$ forward and $\gamma$ backward factors per increment. Also
$\lambda^{*\gamma}_{0,t} = \lambda^{*\gamma}_t$, the transfers
$(\widehat\lambda_t/\widehat\lambda_s)^\gamma$ are continuous in $(s,t)$, and at $\vartheta = \pi$
the transfer of $\lambda^{*\gamma}_{s,t}$ is $\varrho^\gamma$ with $0 < \varrho < 1$, not $1$.

\emph{(2), the subsampling and the computation.} For $\beta \in (0,1)$ write
$g_\beta(m) := N_\beta\beta^{|m|}$ with $N_\beta := \tfrac{1-\beta}{1+\beta}$, and
$\varsigma_\beta := \beta/(1-\beta)^2$. By the transfer above and
$1 - 2\beta\cos\vartheta + \beta^2 = (1-\beta)^2 + 2\beta(1 - \cos\vartheta)$, the transfer of
$g_\beta$ is $\bigl(1 + 2\varsigma_\beta(1 - \cos\vartheta)\bigr)^{-1}$, and $\varsigma_\beta$
increases strictly with $\beta$. Two identities hold for all $m \in \ZZ$:
\[
  g_\beta * g_{\beta'} = \frac{\varsigma_\beta g_\beta - \varsigma_{\beta'}g_{\beta'}}
  {\varsigma_\beta - \varsigma_{\beta'}} \quad (\beta' < \beta), \qquad
  g_\beta(qm) = \frac{N_\beta}{N_{\beta^q}}\,g_{\beta^q}(m) .
\]
The first holds for the transfers, by
$\frac{\varsigma_\beta}{1 + 2\varsigma_\beta u} - \frac{\varsigma_{\beta'}}{1 + 2\varsigma_{\beta'}u}
= \frac{\varsigma_\beta - \varsigma_{\beta'}}{(1 + 2\varsigma_\beta u)(1 + 2\varsigma_{\beta'}u)}$
with $u = 1 - \cos\vartheta$, hence for the sequences. The second is
$\beta^{q|m|} = (\beta^q)^{|m|}$. The kernel of range $T$ on $h\ZZ$ is $g_\beta$ with
$\beta = e^{-h/T}$, so by the second identity its subsample is proportional to
$g_{\beta^q}$, $\beta^q = e^{-qh/T}$, the kernel of range $T$ on $qh\ZZ$. On the grid $q^kh\ZZ$ the
kernel of range $ctq^k$ has the ratio $e^{-q^kh/(ctq^k)} = a$, so it is $g = g_a$ at every level.
By (1) at $s = 0$, where $p = 0$, $g = \lambda_{ct} = \varepsilon^+_{ct} * \varepsilon^-_{ct}$, so
convolution with $g$ is the forward recursion $v(n) = a\,v(n-1) + (1 - a)\,y(n)$ followed by the
same recursion run backward, with coefficients depending on $a$ alone.

Let $F_k$ be $L_{ctq^k}$ sampled on $q^kh\ZZ$ and read on $h\ZZ$: $F_k(i) := L_{ctq^k}(ih)$ when
$q^k$ divides $i$, and $F_k(i) := 0$ otherwise. It is summable, and
$F_k(q^kj) = (2ctq^k)^{-1}a^{|j|} = c_k\,g(j)$ with $c_k := q^{-k}(2ct)^{-1}N_a^{-1} > 0$. Let
$G_n := F_0 * \dots * F_{n-1}$ on $\ZZ$ and $G_0 := \delta_0 = u_0$. By induction
$u_n(m) = \kappa_nG_n(q^nm)$ with $\kappa_0 := 1$ and $\kappa_{n+1} := \kappa_n/c_n$: since $F_n$
vanishes off $q^n\ZZ$,
\[
  G_{n+1}(q^{n+1}m) = \sum_jF_n(q^nj)\,G_n\bigl(q^n(qm - j)\bigr)
  = \frac{c_n}{\kappa_n}\sum_jg(j)\,u_n(qm - j) = \frac{c_n}{\kappa_n}\,u_{n+1}(m) .
\]
This is the restriction statement. For $n = 1$, $u_1(m) = g(qm) = N_aa^{q|m|}$, and
$a^{q|m|} = e^{-qh|m|/(ct)} = 2ct\,L_{ct}(qhm)$, so $U_1$ and $V_1$ are both the renormalized
samples of $L_{ct}$ on $qh\ZZ$.

\emph{(2), the computed level.} Put $x := h/(ct) > 0$, so $a = e^{-x}$, and $\rho_i := a^{q^i}$ for
$i \ge 1$, distinct numbers in $(0, a)$. By induction $u_n(m) = \sum_{i=1}^nb^{(n)}_i\rho_i^{|m|}$
for $n \ge 1$, with $b^{(1)}_1 = N_a$ and
\[
  b^{(n+1)}_{i+1} = b^{(n)}_i\,f_i \ (1 \le i \le n), \qquad
  b^{(n+1)}_1 = \varsigma_aN_a\sum_{i=1}^n\frac{b^{(n)}_i}{N_{\rho_i}(\varsigma_a - \varsigma_{\rho_i})},
  \qquad f_i := -\frac{\varsigma_{\rho_i}}{\varsigma_a - \varsigma_{\rho_i}} .
\]
Write $u_n = \sum_i(b^{(n)}_i/N_{\rho_i})\,g_{\rho_i}$. By the first identity with $\beta = a$ and
$\beta' = \rho_i$, each $g * g_{\rho_i}$ is a combination of $g_a$ and $g_{\rho_i}$, and by the
second, $g_a(qm) = N_a\rho_1^{|m|}$ and $g_{\rho_i}(qm) = N_{\rho_i}\rho_{i+1}^{|m|}$. Collecting
terms gives the recursion. The first rule gives $b^{(n)}_n = N_a\prod_{i=1}^{n-1}f_i$ and, for
$n \ge 2$, $b^{(n)}_{n-1} = b^{(2)}_1\prod_{i=1}^{n-2}f_i$, with
$b^{(2)}_1 = \varsigma_aN_a^2/\bigl(N_{\rho_1}(\varsigma_a - \varsigma_{\rho_1})\bigr)$. Every $f_i$
is nonzero, so both coefficients are, and
$R_n := b^{(n)}_n/b^{(n)}_{n-1} = (N_a/b^{(2)}_1)\,f_{n-1}$ for every $n \ge 2$. For
$\beta = e^{-y}$, $y > 0$, one has $N_\beta = (\cosh y - 1)/\sinh y$ and
$\varsigma_\beta = 1/\bigl(2(\cosh y - 1)\bigr)$, and $\rho_i = e^{-q^ix}$. Substituting,
\[
  R_n = -\frac{\sinh x}{\sinh(qx)}\cdot\frac{\cosh(qx) - \cosh x}{\cosh(q^{n-1}x) - \cosh x} .
\]

\emph{(2), the sampled step.} Put $r_k := ctq^k$, $k \ge 0$, distinct ranges. For $r \ne s$,
\[
  L_r * L_s = \frac{r^2L_r - s^2L_s}{r^2 - s^2}
\]
pointwise. Both sides are even, and for $x \ge 0$ the integral
$\int L_r(y)L_s(x - y)\,dy$, split at $0$ and at $x$, is the sum of three exponential integrals,
\[
  \frac{1}{4rs}\Bigl(\frac{e^{-x/s} + e^{-x/r}}{1/r + 1/s} + \frac{e^{-x/r} - e^{-x/s}}{1/s - 1/r}\Bigr),
\]
which is the right side. By induction on the number of factors,
$L_{r_0} * \dots * L_{r_N} = \sum_{k=0}^NA_kL_{r_k}$: convolving with $L_{r_{N+1}}$ multiplies
$A_k$ by $r_k^2/(r_k^2 - r_{N+1}^2)$ for $k \le N$ and adds a term in $L_{r_{N+1}}$. With
$r_k^2 \propto q^{2k}$ this gives
\[
  A_0 = \prod_{j=1}^{N}\frac{1}{1 - q^{2j}}, \qquad
  A_1 = \frac{q^2}{q^2 - 1}\prod_{j=1}^{N-1}\frac{1}{1 - q^{2j}} \quad (N \ge 1) .
\]
Take $N = n - 1$. At the point $q^nhm$,
$L_{r_k}(q^nhm) = (2r_k)^{-1}e^{-q^nh|m|/(ctq^k)} = (2r_k)^{-1}\rho_{n-k}^{|m|}$. So the samples are
$\sum_{i=1}^n\sigma_i\rho_i^{|m|}$ with $\sigma_i := A_{n-i}/(2r_{n-i})$, and
\[
  \frac{\sigma_n}{\sigma_{n-1}} = \frac{qA_0}{A_1} = -\frac{q^2 - 1}{q\,(q^{2(n-1)} - 1)} .
\]

\emph{(2), the comparison.} Let $n \ge 2$ and suppose that $u_n$ is a real multiple $\kappa$ of
the samples of the step, $u_n(m) = \kappa\sum_i\sigma_i\rho_i^{|m|}$ for every $m$; $U_n = V_n$ is
the case $\kappa > 0$. Then $\sum_{i=1}^n(b^{(n)}_i - \kappa\sigma_i)\rho_i^m = 0$ for every
$m \ge 0$. Geometric sequences with distinct ratios are linearly independent: if
$\sum_{i \in S}d_i\rho_i^m = 0$ for every $m \ge 0$ and $j \in S$, then the shifted sum minus
$\rho_j$ times the sum is $\sum_{i \in S \setminus \{j\}}d_i(\rho_i - \rho_j)\rho_i^m = 0$, so by
induction on the size of $S$ every $d_i(\rho_i - \rho_j)$ with $i \ne j$ vanishes, hence every
$d_i$ with $i \ne j$, and then $d_j$. So $b^{(n)}_i = \kappa\sigma_i$ for every $i$. Since
$b^{(n)}_n \ne 0$, $\kappa \ne 0$, and $R_n = \sigma_n/\sigma_{n-1}$. But
$|R_n| < |\sigma_n/\sigma_{n-1}|$. Put $M := q^{n-1} \ge q$ and $\psi(u) := \sinh u/u$, strictly
increasing on $(0,\infty)$ as $u^2\psi'(u) = u\cosh u - \sinh u$ vanishes at $0$ and has the
derivative $u\sinh u > 0$. First, $\psi(x) < \psi(qx)$ is $\sinh x/\sinh(qx) < 1/q$. Second,
$\cosh(yx) - \cosh x = 2\sinh\bigl((y+1)x/2\bigr)\sinh\bigl((y-1)x/2\bigr)$, so for $y > 1$
\[
  \frac{\cosh(yx) - \cosh x}{y^2 - 1}
  = \frac{x^2}{2}\,\psi\Bigl(\frac{(y+1)x}{2}\Bigr)\,\psi\Bigl(\frac{(y-1)x}{2}\Bigr) ,
\]
nondecreasing in $y$. At $y = q$ and $y = M$ this gives
$0 < (\cosh qx - \cosh x)/(\cosh Mx - \cosh x) \le (q^2 - 1)/(M^2 - 1)$. With the first bound,
$|R_n| < (q^2 - 1)/\bigl(q(M^2 - 1)\bigr) = |\sigma_n/\sigma_{n-1}|$, a contradiction, so
$U_n \ne V_n$.

Every step is machine-checked.
\end{proof}

\paragraph{Measured: the size of the error, and the sampled Gaussian family.} Both ratios of the
proof of (2) tend to the same limit as $x = h/(ct) \to 0$, which is consistent with the error
falling as the grid refines relative to the scale. Its size is measured, not proved. The script
\texttt{scripts/grid\_pyramid\_exactness.py} computes, for $q = 2$, the total-variation distance
between the computed array $U_n$ and the sampled step $V_n$ at the levels $n = 1, \dots, 4$. At
$h/(ct) = 1$ it is $2\cdot10^{-17}$, $6.3\,\%$, $6.7\,\%$ and $6.6\,\%$. At $h/(ct) = 0.5$ it is
$2\cdot10^{-18}$, $2.0\,\%$, $1.3\,\%$ and $1.1\,\%$. At $h/(ct) = 0.25$ it is $1\cdot10^{-16}$,
$0.38\,\%$, $0.25\,\%$ and $0.21\,\%$, and at $h/(ct) = 0.1$ it is $2\cdot10^{-16}$, $0.057\,\%$,
$0.037\,\%$ and $0.033\,\%$. So from level $2$ to level $4$ the computed error stays of one order:
it is flat at $h/(ct) = 1$ and falls by a factor of about two at the finer ratios, covariance holding
$h/(ct)$ fixed. Between ratios it falls fast as $h/(ct)$ does, by about a factor of $100$ from
$h/(ct) = 1$ to $0.1$.

The sampled Gaussian family behaves differently at fine scales. The samples of the Gaussian kernel
on a grid do not form a convolution semigroup [\sscite{lindeberg1990scalespace}, Prop.~12 and the
remark after it, p.~242]. Renormalized, they need not form a hemigroup either. The script
\texttt{scripts/sampled\_gaussian\_increments.py} computes the step between two renormalized
sampled Gaussians as the inverse transform of the ratio of their discrete transfers, on a grid of
$4096$ points. Between the standard deviations $0.5$ and $0.6$ grid units the step has
coefficients as low as $-0.011$, and between $1.0$ and $1.2$ as low as $-0.0046$. So, as computed
at these two pairs of scales, the sampled Gaussian family is not a hemigroup on the grid at the
scales of a pyramid's finest levels, and the negative coefficients shrink as the samples resolve
the kernel. No inequality is proved here for all scales.

On a grid the Laplace family is better behaved than the Gaussian in exactness: the sampled Laplace
family is an exact cascade, and the sampled Gaussian family, as computed above, is not at fine
scales. The reason is the factorization of
Proposition~\ref{prop:sampled-laplace-exact}(1): every step is a convolution of two
one-sided geometric laws mixed with $\delta_0$, with the weight $p \in [0,1)$. It exists because
the sampled Laplace kernel has the transfer of the first display of that proposition, a rational
function of $e^{i\vartheta}$ with one pole inside and one outside the unit circle. The samples of
a Gaussian have a theta function as their transfer instead. The \Matern{} pyramid computed on the
grids is exact in sense (i) at its first level only. The same mask has a second merit,
Proposition~\ref{prop:grid-pyramids-infinitely-divisible}(4): it is infinitely divisible on the
grid, so it generates an exact grid pyramid of its own, in the second stratum. So the grid
offers two different objects: an exact cascade, the sampled Laplace family, and a pyramid
covariant at every level, the geometric pyramid of that family's mask, which is not the sampled
\Matern{} pyramid.

% The section's one relation remark (WRITING section 4; second blind review, L8). It gathers the
% comparisons with the released continuum modules that stood after Example 2.5, as "The continuum
% half" after Proposition 2.21, and in the locality paragraph after Proposition 2.23.
% draft: Remark 2.22 and the paragraph after it (section 2.6), lines 293-297. A comparison with a
% released statement and a reading: prose. The released clause is not quoted here; its statement
% is the interface node thm:cited-A20, read at the page by the librarian. The process references
% of the draft (the risk note on P-0015, P-0001) are dropped.
\paragraph{Relation to the continuum theory.} Three statements of this section have continuum
counterparts in the released modules.

\emph{The generator.} The fully covariant family with the exponent $\Cin(\tau\omega)$, the Cin ray
of step $\tau$ \prog{V-cone v0.1}{Lemma~3.1(1), p.~13}, has at the scale $t$ the generator
$t^{-1}\delta^2_{t\tau}$. This is one symmetric second difference of the generator form
(\ref{eq:generator-form}) \prog{V-cone v0.1}{Prop.~6.4(5), p.~58}
(Theorem~\ref{thm:cited-A13}). Its step grows with the scale. Lindeberg's generator
(Proposition~\ref{prop:lindeberg-stratum-one}) is the same second difference, held at the fixed
step $h$.

\emph{The masks.} The continuum form of
Proposition~\ref{prop:grid-pyramids-infinitely-divisible}(2) is
\prog{V-cone v0.1}{Prop.~7.5(2), p.~65}, stated at the page as Theorem~\ref{thm:cited-A20}. In
paraphrase, for a symmetric rational transform $R$ of an infinitely divisible law it gives three
things. It gives the unique continuous $F$ with $F(0) = 0$ and
$F(q\omega) - F(\omega) = -\log R(\omega)$, the product formula
$e^{-F} = \prod_{m \ge 1}R(q^{-m}\,\cdot)$, and the criterion for $F$ to be an admissible
exponent (\ref{eq:admissible-exponent}). The \Matern{} pyramid is its case
$R(\omega) = (1 + c^2\omega^2)^{-1}$, with $c$ its range factor. The lattice half differs. On
the lattice, rationality of the transfer in $e^{i\vartheta}$ no longer gives infinite
divisibility: the binomial mask has a rational transfer and is not infinitely divisible, so
Proposition~\ref{prop:grid-pyramids-infinitely-divisible}(1) has to supply it. The profile, a
density in the continuum, becomes the finite atomic measure $\varpi_w$, and the interpolant's
profile becomes the atomic $\sum_mT_{q^{-m}}\varpi_w$. That is the same series read on measures,
which is why the product formula survives. The convergence, automatic for a rational $R$, becomes
the logarithmic moment, since a lattice law need have no second moment. The admissibility
criterion of the continuum, a nonincreasing profile, has no lattice counterpart. By
Theorem~\ref{thm:lattice-line}(3) the folded \Levy{} measure of a kernel of a fully covariant
family is absolutely continuous, so an atomic one is never the \Levy{} measure of such a kernel.
So the released clause delivers a fully covariant interpolant, while the pyramid requires the
mask's own covariance under one ratio. The grid meets that weaker demand, and what replaces the
uniqueness of the interpolant is the uniqueness of the knot kernels. Between the knots a grid
pyramid is undetermined. That is the freedom the appendix measures, and here there is no fully
covariant member to select.

\emph{Locality.} On the continuum the Gaussian is the only fully covariant family whose generator
is a differential operator \prog{V-selection v0.1}{Thm.~3.3, p.~16}. The \Matern{} generator is a series
of bounded terms whose sum is unbounded (Proposition~\ref{prop:matern-generator-octave}(1)),
though the \Matern{} family satisfies a differential equation of order three jointly in space and
scale \prog{V-selection v0.1}{Prop.~3.8, p.~24}. Both statements are on the line. So on the continuum the
Gaussian is the better behaved of the two in locality, while on a grid the Laplace family is the
better behaved in exactness.
